% ============================================================================
% RUNNING APPENDIX (progress checklist item E1, assembled 2026-09-29).
% Side and tangent analyses kept outside the paper's main line.  Every number in
% the prose below is a \runapp... macro, and every table a
% generated/appendix_running_<item>.tex, written by
%     python writeup/make_appendix_running_tex.py
% from the committed CSV, study summary or executed-notebook output named in the
% comment above each table.  Nothing here is typed in by hand; rerun the maker
% after any source changes, then rebuild the paper.
% ============================================================================
% AUTO-GENERATED by writeup/make_appendix_running_tex.py -- do not edit.
% Numbers quoted in the prose of writeup/sections/appendix_running.tex; each is
% read from the source named at the table it belongs to (see the maker).
\newcommand{\runappRpNdays}{866}
\newcommand{\runappRpAsMid}{1.90}
\newcommand{\runappRpAsX}{\ensuremath{-}1.92}
\newcommand{\runappRpFlatMid}{3.43}
\newcommand{\runappRpFlatX}{\ensuremath{-}3.62}
\newcommand{\runappRpSignMid}{2.34}
\newcommand{\runappRpSignX}{\ensuremath{-}1.90}
\newcommand{\runappRpHybMid}{2.97}
\newcommand{\runappRpHybX}{\ensuremath{-}1.10}
\newcommand{\runappRpXMin}{\ensuremath{-}3.62}
\newcommand{\runappRpXMax}{\ensuremath{-}1.10}
\newcommand{\runappRpMidMin}{1.90}
\newcommand{\runappRpMidMax}{3.43}
\newcommand{\runappRpXingLo}{16.0}
\newcommand{\runappRpXingHi}{18.2}
\newcommand{\runappRpBeLo}{0.91}
\newcommand{\runappRpBeHi}{1.46}
\newcommand{\runappRpHalfSpread}{1.69}
\newcommand{\runappRpCensored}{13}
\newcommand{\runappRpEightMidLo}{1.76}
\newcommand{\runappRpEightMidHi}{2.47}
\newcommand{\runappRpEightXLo}{\ensuremath{-}2.47}
\newcommand{\runappRpEightXHi}{\ensuremath{-}1.82}
\newcommand{\runappRpRidgeReinv}{2.32}
\newcommand{\runappHcNdays}{866}
\newcommand{\runappHcNbars}{10{,}387}
\newcommand{\runappHcAsMid}{2.93}
\newcommand{\runappHcSignMid}{2.40}
\newcommand{\runappHcAsX}{2.21}
\newcommand{\runappHcSignX}{1.35}
\newcommand{\runappHcSignPctBuy}{38.7}
\newcommand{\runappHcAsLegX}{\ensuremath{-}0.23}
\newcommand{\runappHcSignLegX}{0.86}
\newcommand{\runappCmAccLo}{0.52}
\newcommand{\runappCmAccHi}{0.59}
\newcommand{\runappCmShortLo}{0.61}
\newcommand{\runappCmShortHi}{0.69}
\newcommand{\runappCmBelow}{12}
\newcommand{\runappCmClocks}{12}
\newcommand{\runappCmBuyBeats}{9}
\newcommand{\runappCmSellBeats}{9}
\newcommand{\runappMapNbars}{9{,}521}
\newcommand{\runappMapQStretch}{0.2018}
\newcommand{\runappMapQOne}{0.2093}
\newcommand{\runappMapQTwo}{0.2132}
\newcommand{\runappMapQThree}{0.2009}
\newcommand{\runappMapAsMid}{3.38}
\newcommand{\runappMapAsX}{2.65}
\newcommand{\runappMapSmMid}{2.30}
\newcommand{\runappMapSmX}{1.25}
\newcommand{\runappMapOneMid}{1.70}
\newcommand{\runappMapOneX}{0.68}
\newcommand{\runappMapOneDis}{10.4}
\newcommand{\runappMapOneBuy}{46.6}
\newcommand{\runappMapTwoMid}{1.15}
\newcommand{\runappMapTwoX}{0.22}
\newcommand{\runappMapTwoDis}{15.9}
\newcommand{\runappMapTwoBuy}{48.3}
\newcommand{\runappMapThreeMid}{0.86}
\newcommand{\runappMapThreeX}{\ensuremath{-}0.07}
\newcommand{\runappMapThreeDis}{17.3}
\newcommand{\runappMapThreeBuy}{48.4}
\newcommand{\runappMapSmBuy}{38.6}
\newcommand{\runappMapOneDs}{\ensuremath{-}0.60}
\newcommand{\runappMapOneCi}{[\ensuremath{-}1.11, \ensuremath{-}0.11]}
\newcommand{\runappMapExcl}{6}
\newcommand{\runappMapIntervals}{6}
\newcommand{\runappPcAsPeak}{4.20}
\newcommand{\runappPcAsPeakClock}{11:00}
\newcommand{\runappPcAsSecond}{4.18}
\newcommand{\runappPcAsSecondClock}{10:30}
\newcommand{\runappPcAsFirst}{3.11}
\newcommand{\runappPcAsFirstX}{2.44}
\newcommand{\runappPcSgFirst}{0.68}
\newcommand{\runappPcSgFirstX}{0.02}
\newcommand{\runappPcAsFifteen}{1.93}
\newcommand{\runappPcAsLast}{0.24}
\newcommand{\runappPcAsLastX}{\ensuremath{-}0.23}
\newcommand{\runappPcSgLast}{1.30}
\newcommand{\runappPcSgLastX}{0.83}
\newcommand{\runappPcSgLastLo}{0.92}
\newcommand{\runappPcSgLastHi}{1.51}
\newcommand{\runappPcDayClocks}{12}
\newcommand{\runappPcAsBeatsAll}{12}
\newcommand{\runappPcAsBeatsAllX}{12}
\newcommand{\runappPcSgDayLo}{0.67}
\newcommand{\runappPcSgDayHi}{1.97}
\newcommand{\runappPcRpClocks}{12}
\newcommand{\runappPcRpAsLo}{0.20}
\newcommand{\runappPcRpAsHi}{2.51}
\newcommand{\runappPcRpSgBeats}{5}
\newcommand{\runappPcNbDiff}{0.06}
\newcommand{\runappPcNbLast}{1.33}
\newcommand{\runappRdBaseLo}{\ensuremath{-}19}
\newcommand{\runappRdBaseHi}{\ensuremath{-}18}
\newcommand{\runappRdLiveLo}{\ensuremath{-}37}
\newcommand{\runappRdLiveHi}{\ensuremath{-}31}
\newcommand{\runappRdAllLo}{\ensuremath{-}44}
\newcommand{\runappRdAllHi}{\ensuremath{-}36}
\newcommand{\runappRdBetterLo}{6}
\newcommand{\runappRdBetterHi}{11}
\newcommand{\runappRdNlo}{1{,}404}
\newcommand{\runappRdNhi}{1{,}406}
\newcommand{\runappRdFirst}{2018-09-24}
\newcommand{\runappRdLast}{2024-04-30}
\newcommand{\runappRdLastMax}{0.018}
\newcommand{\runappRdAlpha}{0.05}
\newcommand{\runappRdExtremeN}{5}
\newcommand{\runappRdCells}{108}
\newcommand{\runappRdExtremeQ}{132.0}
\newcommand{\runappRdExtremeOne}{0.197}
\newcommand{\runappRdExtremeClock}{14:00}
\newcommand{\runappRdExtremeCell}{all features ridge}
\newcommand{\runappRdAllRidgeLast}{2.65}
\newcommand{\runappPoEightyNineResDays}{814}
\newcommand{\runappPoEightySevenDays}{814}
\newcommand{\runappPoSpxDays}{395}
\newcommand{\runappPoSpxBuys}{212}
\newcommand{\runappPoSpxPnl}{\ensuremath{-}0.0129}
\newcommand{\runappPoSpxT}{\ensuremath{-}0.31}
\newcommand{\runappPoResPnl}{0.0538}
\newcommand{\runappPoResT}{1.60}
\newcommand{\runappPoResDays}{564}
\newcommand{\runappPoPostPnl}{\ensuremath{-}0.0114}
\newcommand{\runappPoPostT}{\ensuremath{-}0.35}
\newcommand{\runappPoPostDays}{568}
\newcommand{\runappPoPostCi}{[\ensuremath{-}0.075, 0.057]}
\newcommand{\runappPoResShare}{38}
\newcommand{\runappPoPostShare}{51}
\newcommand{\runappPoDeckPnl}{0.0455}
\newcommand{\runappPoAllPnl}{0.0360}
\newcommand{\runappPoAllD}{\ensuremath{-}0.0095}
\newcommand{\runappPoAllCi}{[\ensuremath{-}0.045, 0.022]}
\newcommand{\runappPoAloneMax}{0.014}
\newcommand{\runappPoFOnePost}{0.0225}
\newcommand{\runappPoFOneT}{0.77}
\newcommand{\runappPoFOneCi}{[\ensuremath{-}0.033, 0.084]}
\newcommand{\runappPoFOneD}{0.034}
\newcommand{\runappPoFOneDCi}{[0.009, 0.061]}
\newcommand{\runappPoFOneHolm}{0.014}
\newcommand{\runappPoFOneResD}{\ensuremath{-}0.033}
\newcommand{\runappPoFOneResDCi}{[\ensuremath{-}0.078, \ensuremath{-}0.003]}
\newcommand{\runappPoTwLo}{\ensuremath{-}0.024}
\newcommand{\runappPoTwHi}{\ensuremath{-}0.012}
\newcommand{\runappPoRatioShareLo}{2}
\newcommand{\runappPoRatioShareHi}{18}
\newcommand{\runappPoRatioPnlLo}{\ensuremath{-}0.009}
\newcommand{\runappPoRatioPnlHi}{0.002}
\newcommand{\runappPoQCard}{0.185}
\newcommand{\runappPoQAsk}{0.142}
\newcommand{\runappPoQModelLo}{0.185}
\newcommand{\runappPoQModelHi}{0.193}
\newcommand{\runappPoQCardRes}{0.097}
\newcommand{\runappPoQAskRes}{0.110}
\newcommand{\runappPoCFourPost}{0.0073}
\newcommand{\runappPoCFourCi}{[\ensuremath{-}0.021, 0.058]}
\newcommand{\runappDtDays}{866}
\newcommand{\runappDtOneR}{0.152}
\newcommand{\runappDtOneT}{1.97}
\newcommand{\runappDtOneN}{286}
\newcommand{\runappDtTwoR}{0.010}
\newcommand{\runappDtTwoT}{0.17}
\newcommand{\runappDtTwoD}{\ensuremath{-}0.046}
\newcommand{\runappDtTwoCi}{[\ensuremath{-}0.082, \ensuremath{-}0.016]}
\newcommand{\runappDsDays}{403}
\newcommand{\runappDsMonthEnds}{40}
\newcommand{\runappDsHolmMin}{0.773}
\newcommand{\runappDsTmax}{1.48}
\newcommand{\runappDsMeLast}{0.356}
\newcommand{\runappDsMeLastT}{1.45}
\newcommand{\runappDsMeLaterLo}{0.11}
\newcommand{\runappDsMeLaterHi}{0.13}
\newcommand{\runappDsMeHolmMin}{1.00}
\newcommand{\runappDsSpreadLo}{2.7}
\newcommand{\runappDsSpreadHi}{5.7}
\newcommand{\runappDsXspSpreadLo}{9.8}
\newcommand{\runappDsXspSpreadHi}{40.0}
\newcommand{\runappDsFirst}{2024-12-06}
\newcommand{\runappDsLast}{2026-09-21}
\newcommand{\runappDgFineLo}{\ensuremath{-}2.2}
\newcommand{\runappDgFineHi}{0.2}
\newcommand{\runappDgFineDm}{0.66}
\newcommand{\runappDgDays}{445}
\newcommand{\runappDgLong}{\ensuremath{-}1.9}
\newcommand{\runappLoDays}{259}
\newcommand{\runappLoOnce}{0.0937}
\newcommand{\runappLoOnceT}{1.70}
\newcommand{\runappLoOneMin}{\ensuremath{-}0.040}
\newcommand{\runappLoOneMinT}{\ensuremath{-}2.85}
\newcommand{\runappMeN}{70}
\newcommand{\runappMeFirst}{2020-01-31}
\newcommand{\runappMeLast}{2025-12-31}
\newcommand{\runappMeR}{0.432}
\newcommand{\runappMeT}{2.62}
\newcommand{\runappMeFirstHalf}{0.189}
\newcommand{\runappMeSecondHalf}{0.676}
\newcommand{\runappMeSplit}{2022-12-30}
\newcommand{\runappMeEarlyLo}{0.152}
\newcommand{\runappMeEarlyHi}{0.422}
\newcommand{\runappMeHolmMin}{1.00}
\newcommand{\runappVxQ}{\ensuremath{-}5.9}
\newcommand{\runappVxQThirteen}{\ensuremath{+}3.3}
\newcommand{\runappVxMid}{1.48}
\newcommand{\runappVxX}{1.00}
\newcommand{\runappVxLfMid}{1.90}
\newcommand{\runappVxLfX}{1.44}
\newcommand{\runappVxNbMid}{1.38}
\newcommand{\runappVxNbLfMid}{1.68}
\newcommand{\runappVxMinusQ}{\ensuremath{-}0.04}
\newcommand{\runappVxFreeQ}{\ensuremath{-}0.43}
\newcommand{\runappVxFreeMid}{1.83}
\newcommand{\runappVxDays}{866}
\newcommand{\runappKeyLive}{\texttt{sumret}, \texttt{sumabsret}, \texttt{sumret3}, \texttt{sumret4}, \texttt{sumpret2}, \texttt{sumbipow}, \texttt{sumautocov}, \texttt{sumvolume}, \texttt{numobs}, \texttt{vix}, \texttt{vvix}, \texttt{vix3m}, \texttt{fomc\_release}, \texttt{fomc\_day}, \texttt{fomc\_until\_inv}, \texttt{fomc\_since\_inv}}
\newcommand{\runappKeyFreeDrop}{\texttt{numobs}, \texttt{vvix}, \texttt{vix3m}}
\newcommand{\runappPbSessLo}{\ensuremath{-}9.0}
\newcommand{\runappPbSessHi}{\ensuremath{-}3.3}
\newcommand{\runappPbSessCi}{3}
\newcommand{\runappPbN}{7}
\newcommand{\runappPbFomcBarLo}{\ensuremath{-}39.0}
\newcommand{\runappPbFomcBarHi}{\ensuremath{-}30.5}
\newcommand{\runappPbDays}{866}
\newcommand{\runappPbSixteenWho}{ridge (HAR + calendar)}
\newcommand{\runappPbSixteen}{\ensuremath{-}6.3}
\newcommand{\runappPbSixteenCi}{[\ensuremath{-}10.4, \ensuremath{-}1.9]}
\newcommand{\runappPbFomcDays}{33}
\newcommand{\runappPbOtherDays}{833}
\newcommand{\runappPbFomcLo}{\ensuremath{-}78}
\newcommand{\runappPbFomcHi}{\ensuremath{-}43}
\newcommand{\runappPbOtherLo}{\ensuremath{-}16}
\newcommand{\runappPbOtherHi}{\ensuremath{-}5}
\newcommand{\runappPbShareLo}{0.45}
\newcommand{\runappPbShareHi}{0.54}
\newcommand{\runappMhDays}{5{,}582}
\newcommand{\runappMhFirst}{2002-01-23}
\newcommand{\runappMhLast}{2024-04-30}
\newcommand{\runappMhPostDays}{2{,}089}
\newcommand{\runappMhNaive}{0.863}
\newcommand{\runappMhAOne}{0.1600}
\newcommand{\runappMhBOne}{0.1435}
\newcommand{\runappMhBThree}{0.1367}
\newcommand{\runappMhDmLo}{3.2}
\newcommand{\runappMhDmHi}{7.2}
\newcommand{\runappMhFourteen}{0.196}
\newcommand{\runappMhFourteenThirty}{0.122}
\newcommand{\runappMhMissTen}{9.4}
\newcommand{\runappMhMissFifteen}{54.6}
\newcommand{\runappSzForecasts}{15}
\newcommand{\runappSzRankAbove}{8}
\newcommand{\runappSzRankMed}{\ensuremath{+}0.03}
\newcommand{\runappSzSignMid}{1.54}
\newcommand{\runappSzDays}{614}
\newcommand{\runappPwDays}{614}
\newcommand{\runappPwFirst}{2021-08-11}
\newcommand{\runappPwSpx}{0.40}
\newcommand{\runappPwComb}{1.44}
\newcommand{\runappPwGain}{1.04}
\newcommand{\runappPwGainCi}{[\ensuremath{-}0.46, 2.55]}
\newcommand{\runappPwShare}{0.73}
\newcommand{\runappPwDollar}{4.32}
\newcommand{\runappPwCorrLo}{\ensuremath{-}0.04}
\newcommand{\runappPwCorrHi}{\ensuremath{+}0.03}
\newcommand{\runappPwAbsLo}{\ensuremath{-}0.23}
\newcommand{\runappPwAbsHi}{\ensuremath{-}0.08}
\newcommand{\runappPwBestLo}{\ensuremath{-}0.11}
\newcommand{\runappPwForecasts}{15}
\newcommand{\runappCoExA}{1}
\newcommand{\runappCoExB}{0}
\newcommand{\runappCoSeries}{13}
\newcommand{\runappCoOne}{\ensuremath{-}0.11}
\newcommand{\runappCoOneCi}{[\ensuremath{-}0.206, \ensuremath{-}0.001]}


% Shrink a table to the text width only when it is wider; never enlarge it.
\newsavebox{\runappbox}
\newcommand{\runappfit}[1]{%
  \sbox{\runappbox}{#1}%
  \ifdim\wd\runappbox>\textwidth
    \resizebox{\textwidth}{!}{\usebox{\runappbox}}%
  \else
    \usebox{\runappbox}%
  \fi}

\section{Running Appendix: Side Analyses}
\label{sec:appendix_running}

This appendix collects analyses run alongside the paper that are not part of
its argument: other entry clocks and holding periods for the same-day option
trade, other forecast targets, the trade after the forecast panel ends, and
several checks on the last-30-min trade itself. Each subsection states the
question, gives the table, states the finding, and says why it is kept out of
the main text.

\paragraph{Conventions.} Unless a subsection says otherwise: the instrument is
the straddle of Section~\ref{sec:close_option_methods} (the nearest
out-of-the-money call and put of the same-day expiry, one position); $R$ is the
return of the long straddle per unit of midpoint premium and a short earns
$-R$; Sharpe ratios are annualized with $\sqrt{252}$ per trade day, as in
Table~\ref{tab:close_option_rules}; a crossed fill buys at the ask and sells at
the bid at entry, and a position that settles in cash pays no exit spread; the
forecast is the block-diagonal ridge; and the scored frame is the
\runappHcNdays{} expiration days of Table~\ref{tab:close_option_rules}. The
last-30-min trade is the trade of Section~\ref{sec:close_option_results}: entry
at 15:30, cash settlement at the close. Panel stamps are bar-end labelled, so
the forecast used at a clock $t$ is the one issued at $t$ for the bar
$[t, t+30]$.

% ----------------------------------------------------------------------------
\subsection{Every Thirty Minutes: A Re-Picked Straddle}
\label{sec:app_run_repick}

\paragraph{Question.} Does the forecast trade at every thirty-minute bar of the
session, not only the last? At each clock $t$ from 10:00 to 15:00 the nearest
out-of-the-money straddle is bought or sold at its midpoint and closed at the
next stamp's midpoint of the same strikes, a thirty-minute hold; the 15:30
straddle is held to settlement, so that leg is the last-30-min trade. The
strikes are picked afresh at every clock. The chain carries only the same-day
expiry, so no thirty-minute implied variance is quoted. It is built from the
quoted variance to the close, $\mathrm{IV}^2_{\mathrm{hr},t}h_t$ (the vendor's
hourly implied volatility squared, times the hours left), multiplied by $w_t$,
the bar's average share of the rest of the day's realized variance on earlier
days. The signal is
$s^{\mathrm m}_t=\widehat{RV}_t-\mathrm{IV}^2_{\mathrm{hr},t}\,h_t\,w_t$, and a
day's return is the sum of its twelve legs.
Table~\ref{tab:app_run_repick} scores four rules: always short; always short
with the 15:30 leg left flat; $\mathrm{sign}(s^{\mathrm m})$ at every clock;
and always short before 15:30 with $\mathrm{sign}(s^{\mathrm m})$ on the 15:30
leg.

% Source: results/atm_straddle_intraday/rule_table_intraday_blk2.csv and
%   rule_table_intraday_crossed_blk2.csv (notebooks/atm_straddle_intraday.executed.ipynb,
%   sections 6-7), panel (b) results/atm_straddle_intraday/rule_by_strategy/pooled/
%   rule_by_strategy_sign_s.csv (writeup/make_rule_by_strategy_intraday_tex.py, the bundle
%   writeup/rule_by_strategy_intraday_index.pdf).  That bundle was not rebuilt with the
%   2026-09-18 chain's 09:35 stamp; the re-pick's scored frame starts at 10:00 either way.
\begin{table}[H]
\centering
\caption{The re-picked straddle, pooled over the twelve clocks 10:00--15:30 on
the \runappRpNdays{} expiration days. mean/day and $t$ belong to the daily sum
of the twelve legs, and the Sharpe ratio is that sum's mean over its standard
deviation times $\sqrt{252}$. Crossings per day count the trades that cross
the spread (a re-pick onto the same strikes with the same sign is a hold); the
break-even half-spread is the mean daily midpoint profit over the mean
crossings per day, in percent of premium. Panel (a) leaves flat the
\runappRpCensored{} bars whose vendor implied volatility is a censored solver
value; panel (b) re-inverts it from the straddle midpoint, as the main text
does, which moves the block-diagonal ridge from \runappRpSignMid{} to
\runappRpRidgeReinv{}.}
\label{tab:app_run_repick}
\runappfit{\small% AUTO-GENERATED by writeup/make_appendix_running_tex.py -- do not edit.
% Source: results/atm_straddle_intraday/rule_table_intraday_blk2.csv
% Source: results/atm_straddle_intraday/rule_table_intraday_crossed_blk2.csv
% Source: results/atm_straddle_intraday/rule_by_strategy/pooled/rule_by_strategy_sign_s.csv
\begin{tabular}{lrrrrrrr}
\toprule
 & mean/day & $t$ & \shortstack{Sharpe\\(mid)} & \shortstack{Sharpe\\(crossed)} & \shortstack{crossings\\per day} & \shortstack{break-even half-\\spread (\% prem.)} & \shortstack{\% bars\\bought} \\
\midrule
\multicolumn{8}{l}{\emph{(a) four rules, block-diagonal ridge}} \\
always short & 0.160 & 3.53 & 1.90 & \ensuremath{-}1.92 & 16.4 & 0.98 & 0.0 \\
always short, flat at 15:30 & 0.146 & 6.36 & 3.43 & \ensuremath{-}3.62 & 16.0 & 0.91 & 0.0 \\
$\mathrm{sign}(s^{\mathrm m})$ & 0.200 & 4.35 & 2.34 & \ensuremath{-}1.90 & 18.2 & 1.10 & 38.7 \\
always short, $\mathrm{sign}(s^{\mathrm m})$ at 15:30 & 0.243 & 5.51 & 2.97 & \ensuremath{-}1.10 & 16.7 & 1.46 & 3.3 \\
\midrule
\multicolumn{8}{l}{\emph{(b) $\mathrm{sign}(s^{\mathrm m})$ for each forecast}} \\
baseline (HAR + calendar OLS) & 0.163 & 3.54 & 1.91 & \ensuremath{-}2.35 & --- & --- & --- \\
block-diagonal ridge & 0.198 & 4.30 & 2.32 & \ensuremath{-}1.94 & --- & --- & --- \\
block-diagonal ridge, without the FOMC columns & 0.192 & 4.19 & 2.26 & \ensuremath{-}2.02 & --- & --- & --- \\
LightGBM & 0.179 & 3.92 & 2.12 & \ensuremath{-}2.16 & --- & --- & --- \\
XGBoost & 0.173 & 3.75 & 2.02 & \ensuremath{-}2.21 & --- & --- & --- \\
lasso (causally tuned) & 0.180 & 3.96 & 2.13 & \ensuremath{-}2.17 & --- & --- & --- \\
lasso ($\alpha=10^{-4}$) & 0.209 & 4.59 & 2.47 & \ensuremath{-}1.82 & --- & --- & --- \\
elastic net (causally tuned) & 0.151 & 3.26 & 1.76 & \ensuremath{-}2.47 & --- & --- & --- \\
\bottomrule
\end{tabular}
}
\end{table}

\paragraph{Finding.} At the midpoint every rule earns (Sharpe
\runappRpMidMin{} to \runappRpMidMax{}) and $\mathrm{sign}(s^{\mathrm m})$ is
ahead of always short (\runappRpSignMid{} against \runappRpAsMid{}; the eight
forecasts \runappRpEightMidLo{} to \runappRpEightMidHi{}). At the crossed
spread every rule loses (\runappRpXMin{} to \runappRpXMax{}; the eight
forecasts \runappRpEightXLo{} to \runappRpEightXHi{}). The book crosses the
spread \runappRpXingLo{} to \runappRpXingHi{} times a day, and the median
half-spread, \runappRpHalfSpread{}\% of premium, is above every rule's
break-even half-spread (\runappRpBeLo{}\% to \runappRpBeHi{}\%). The least
negative rule at the crossed spread is the one that uses the forecast only on
the settling 15:30 leg (\runappRpHybX{}).

\paragraph{Why it is not in the main text.} No rule survives the spread, and
the implied side of the daytime signal is a construction --- a profile share of
a quote for the rest of the day --- not a price.

% ----------------------------------------------------------------------------
\subsection{Time to the Close: Holding from Every Clock, Delta-Hedged}
\label{sec:app_run_holdclose}

\paragraph{Question.} Enter at any clock and hold to the close: does a
forecast of the rest of the day, read against the quoted variance for the rest
of the day, pick the side? At each clock $t$ from 10:00 to 15:30 the nearest
out-of-the-money straddle is bought or sold at its midpoint and held to cash
settlement, twelve overlapping trades a day. Every position is delta-hedged in
the index: at the entry stamp and at every later thirty-minute stamp to 15:30,
the straddle's Black--Scholes delta (entry strikes, that stamp's spot, and that
stamp's implied volatility over the hours left) is offset by index units held
to the next stamp, the last from 15:30 to the official close. The hedge trades
at the vendor spot at no cost, and $R_t$ adds its profit to the settlement
value. The signal stretches the forecast to the rest of the day,
$s_t=\widehat{RV}_t/w_t-\mathrm{IV}^2_{\mathrm{hr},t}h_t$, which has the sign
of $s^{\mathrm m}_t$ above; at 15:30, $w_t=1$ and $h_t=\tfrac12$, and $s_t$ is
the signal of the last-30-min trade.

% Source: results/atm_straddle_intraday_holdclose/rule_table_intraday_blk2.csv and
%   rule_table_intraday_crossed_blk2.csv (notebooks/atm_straddle_intraday_holdclose.executed.ipynb,
%   sections 4b, 6 and 7; the hedge is attach_long_dh of
%   writeup/make_rule_by_strategy_intraday_tex.py, gated in section 4b).
\begin{table}[H]
\centering
\caption{Holding to the close from every clock, delta-hedged, pooled over the
twelve entry clocks 10:00--15:30 (\runappHcNbars{} trades); columns as in
Table~\ref{tab:app_run_repick}. Each trade crosses the spread once, at entry.
The 15:30 leg column scores that leg alone at the crossed spread.}
\label{tab:app_run_holdclose}
\runappfit{\small% AUTO-GENERATED by writeup/make_appendix_running_tex.py -- do not edit.
% Source: results/atm_straddle_intraday_holdclose/rule_table_intraday_blk2.csv
% Source: results/atm_straddle_intraday_holdclose/rule_table_intraday_crossed_blk2.csv
\begin{tabular}{lrrrrrrrr}
\toprule
 & days & mean/day & $t$ & \shortstack{Sharpe\\(mid)} & \shortstack{Sharpe\\(crossed)} & \shortstack{mean/day\\(crossed)} & \shortstack{15:30 leg Sharpe\\(crossed)} & \shortstack{\% bars\\bought} \\
\midrule
always short & 866 & 1.014 & 5.44 & 2.93 & 2.21 & 0.766 & \ensuremath{-}0.23 & 0.0 \\
$\mathrm{sign}(s)$ & 866 & 0.563 & 4.46 & 2.40 & 1.35 & 0.315 & 0.86 & 38.7 \\
\bottomrule
\end{tabular}
}
\end{table}

% Source: results/atm_straddle_intraday_holdclose/confusion_matrix_by_clock.csv
%   (holdclose notebook, section 8c).
\begin{table}[H]
\centering
\caption{The position against the outcome, by entry clock (midpoint, hedged
$R$). Accuracy is the share of days on which the position had the sign of
$R$: a buy with $R>0$ or a sale with $R\le 0$. The accuracy of always short is
$1-P(R>0)$. The bars with a censored implied volatility are left out.}
\label{tab:app_run_confusion}
\runappfit{\small% AUTO-GENERATED by writeup/make_appendix_running_tex.py -- do not edit.
% Source: results/atm_straddle_intraday_holdclose/confusion_matrix_by_clock.csv
\begin{tabular}{lrrrrrrr}
\toprule
entry & days & \shortstack{accuracy\\$\mathrm{sign}(s)$} & \shortstack{accuracy\\always short} & \shortstack{buy\\share} & $P(R>0)$ & $P(R>0\mid\text{buy})$ & $P(R\le 0\mid\text{sell})$ \\
\midrule
10:00 & 864 & 0.573 & 0.660 & 0.351 & 0.340 & 0.376 & 0.679 \\
10:30 & 865 & 0.566 & 0.686 & 0.348 & 0.314 & 0.329 & 0.693 \\
11:00 & 865 & 0.525 & 0.679 & 0.422 & 0.321 & 0.318 & 0.676 \\
11:30 & 864 & 0.565 & 0.675 & 0.388 & 0.325 & 0.358 & 0.696 \\
12:00 & 866 & 0.539 & 0.630 & 0.424 & 0.370 & 0.392 & 0.647 \\
12:30 & 866 & 0.540 & 0.648 & 0.426 & 0.352 & 0.374 & 0.664 \\
13:00 & 866 & 0.545 & 0.642 & 0.483 & 0.358 & 0.400 & 0.681 \\
13:30 & 866 & 0.585 & 0.648 & 0.312 & 0.352 & 0.400 & 0.669 \\
14:00 & 864 & 0.517 & 0.634 & 0.492 & 0.366 & 0.381 & 0.649 \\
14:30 & 865 & 0.518 & 0.609 & 0.394 & 0.391 & 0.384 & 0.605 \\
15:00 & 862 & 0.566 & 0.618 & 0.215 & 0.382 & 0.378 & 0.617 \\
15:30 & 861 & 0.551 & 0.627 & 0.400 & 0.373 & 0.404 & 0.648 \\
\bottomrule
\end{tabular}
}
\end{table}

\paragraph{A forecast fitted on the rest of the day.} The stretch
$\widehat{RV}_t/w_t$ reaches the rest of the day from a one-bar forecast. Three
maps fitted on the rest-of-day variance itself replace it, one regression per
entry clock for $\log RV^{\mathrm{rem}}_t$ on sessions strictly before the day:
M1 on the log of the next-bar forecast, M2 adding the realized variance from
the open to $t$, and M3 adding the previous session's rest-of-day variance
(the constructions of Section~\ref{sec:app_run_multihorizon}). The 15:30 clock
is left out because there the rest of the day is one bar.

% Source: results/atm_straddle_intraday_holdclose/rule_table_srem_blk2.csv and
%   srem_vs_sm_dsharpe.csv; the QLIKE column from the saved output of the holdclose
%   notebook, section 8d (pooled line of its rest-of-day QLIKE table).
\begin{table}[H]
\centering
\caption{Rest-of-day signals on the eleven entry clocks 10:00--15:00
(\runappMapNbars{} trades, the \runappHcNdays{} days). QLIKE scores each
forecast against the realized variance from $t$ to the close.
$\mathrm{sign}(s^{\mathrm{rem}})$ trades $F^{\mathrm{rem}}_t-
\mathrm{IV}^2_{\mathrm{hr},t}h_t$ with the fitted map in place of the stretch.
$\Delta$Sharpe is the paired difference from the profile-stretch row at the
midpoint, with a 95\% moving-block bootstrap percentile interval.}
\label{tab:app_run_maps}
\runappfit{\small% AUTO-GENERATED by writeup/make_appendix_running_tex.py -- do not edit.
% Source: results/atm_straddle_intraday_holdclose/rule_table_srem_blk2.csv
% Source: results/atm_straddle_intraday_holdclose/srem_vs_sm_dsharpe.csv
% Source: notebooks/atm_straddle_intraday_holdclose.executed.ipynb (section 8d, pooled QLIKE line of the saved output)
\begin{tabular}{lrrrrrrr}
\toprule
 & \shortstack{QLIKE,\\rest of day} & \shortstack{Sharpe\\(mid)} & \shortstack{Sharpe\\(crossed)} & \shortstack{\% bars\\bought} & \shortstack{\% bars sign\\differs} & \shortstack{$\Delta$Sharpe\\(mid)} & \shortstack{95\%\\interval} \\
\midrule
always short (no forecast) & --- & 3.38 & 2.65 & 0.0 & --- & --- & --- \\
$\mathrm{sign}(s)$, profile stretch $\widehat{RV}_t/w_t$ & 0.2018 & 2.30 & 1.25 & 38.6 & --- & --- & --- \\
$\mathrm{sign}(s^{\mathrm{rem}})$, M1 & 0.2093 & 1.70 & 0.68 & 46.6 & 10.4 & \ensuremath{-}0.60 & [\ensuremath{-}1.11, \ensuremath{-}0.11] \\
$\mathrm{sign}(s^{\mathrm{rem}})$, M2 & 0.2132 & 1.15 & 0.22 & 48.3 & 15.9 & \ensuremath{-}1.15 & [\ensuremath{-}1.82, \ensuremath{-}0.59] \\
$\mathrm{sign}(s^{\mathrm{rem}})$, M3 & 0.2009 & 0.86 & \ensuremath{-}0.07 & 48.4 & 17.3 & \ensuremath{-}1.44 & [\ensuremath{-}2.14, \ensuremath{-}0.85] \\
\bottomrule
\end{tabular}
}
\end{table}

\paragraph{Finding.} Held to the close and hedged, both rules pay their spread
(Table~\ref{tab:app_run_holdclose}): always short \runappHcAsMid{} at the
midpoint and \runappHcAsX{} crossed, $\mathrm{sign}(s)$ \runappHcSignMid{} and
\runappHcSignX{}. The forecast-free control is ahead at both fills. The
direction rule is right on \runappCmAccLo{} to \runappCmAccHi{} of days at each
clock; always selling is right on \runappCmShortLo{} to \runappCmShortHi{}, more
often at \runappCmBelow{} of the \runappCmClocks{} clocks
(Table~\ref{tab:app_run_confusion}). Where it buys, the straddle wins more
often than on an average day at \runappCmBuyBeats{} of the \runappCmClocks{}
clocks. The fitted rest-of-day maps forecast the rest of the day no better than
the stretch (pooled QLIKE \runappMapQOne{}, \runappMapQTwo{} and
\runappMapQThree{} against \runappMapQStretch{}), buy more often
(\runappMapOneBuy{}\% of bars for M1 against \runappMapSmBuy{}\%), and trade
worse: Sharpe \runappMapOneMid{}, \runappMapTwoMid{} and \runappMapThreeMid{}
against \runappMapSmMid{} at the midpoint, with all \runappMapExcl{} of
\runappMapIntervals{} paired intervals (two fills, three maps) below zero; M1's
difference is \runappMapOneDs{} \runappMapOneCi{}.

\paragraph{Why it is not in the main text.} Before 15:30 the forecast does not
improve on the forecast-free control, the daytime trades rest on a
frictionless hedge the main text does not model, and the twelve trades of a day
overlap.

% ----------------------------------------------------------------------------
\subsection{The Sharpe Ratio Through the Day}
\label{sec:app_run_perclock}

\paragraph{Question.} How does the Sharpe ratio of each rule change with the
clock at which the position is taken? Table~\ref{tab:app_run_perclock} keeps
one clock at a time, one trade per expiration day, for the two books above:
the thirty-minute re-pick (unhedged) and the delta-hedged hold to the close.
The 09:35 row is the first stamp of the chain; it is scored with the forecast
issued at 09:30 for the 09:30--10:00 bar, and only in the hold-to-close book.

% Source: results/atm_straddle_intraday/rule_by_strategy/<hhmm>/rule_by_strategy_{always_short,sign_s}.csv
%   and results/atm_straddle_intraday_holdclose/rule_by_strategy_dh/<hhmm>/rule_by_strategy_{always_short,sign_s}.csv
%   (writeup/make_rule_by_strategy_intraday_tex.py [--dh-holdclose]; the bundles
%   writeup/rule_by_strategy_intraday_index.pdf and writeup/rule_by_strategy_dh_holdclose_index.pdf,
%   the latter current, the former not rebuilt with 09:35).  These re-invert the censored
%   implied volatility as the main text does; the notebooks' own per-clock tables leave those
%   bars flat (caption).
\begin{table}[H]
\centering
\caption{Annualized Sharpe ratio by entry clock, one trade per expiration day.
Re-pick: the thirty-minute hold of Section~\ref{sec:app_run_repick} (its 15:30
row is the last-30-min trade). To close: the delta-hedged hold of
Section~\ref{sec:app_run_holdclose}. $\mathrm{sign}(s)$ is the block-diagonal
ridge; the last column is the range over the eight forecasts of
Table~\ref{tab:close_option_rules}. Midpoint fills unless marked crossed. The
censored implied volatility is re-inverted as in the main text; the
hold-to-close notebook, which leaves those bars flat, differs from the
$\mathrm{sign}(s)$ column by at most \runappPcNbDiff{} (15:30:
\runappPcNbLast{}).}
\label{tab:app_run_perclock}
\runappfit{\small% AUTO-GENERATED by writeup/make_appendix_running_tex.py -- do not edit.
% Source: results/atm_straddle_intraday/rule_by_strategy/<hhmm>/rule_by_strategy_{always_short,sign_s}.csv (30-minute re-pick)
% Source: results/atm_straddle_intraday_holdclose/rule_by_strategy_dh/<hhmm>/rule_by_strategy_{always_short,sign_s}.csv (hold to close, delta-hedged)
\begin{tabular}{lrrrrrrrr}
\toprule
entry & days & \shortstack{re-pick\\always short} & \shortstack{re-pick\\$\mathrm{sign}(s)$} & \shortstack{to close\\always short} & \shortstack{to close\\always short, crossed} & \shortstack{to close\\$\mathrm{sign}(s)$} & \shortstack{to close\\$\mathrm{sign}(s)$, crossed} & \shortstack{to close, $\mathrm{sign}(s)$\\eight forecasts} \\
\midrule
09:35 & 863 & --- & --- & 3.11 & 2.44 & 0.68 & 0.02 & 0.62 to 1.71 \\
10:00 & 864 & 0.26 & \ensuremath{-}0.11 & 3.32 & 2.29 & 1.66 & 0.65 & 1.32 to 2.04 \\
10:30 & 865 & 0.64 & 0.39 & 4.18 & 3.51 & 1.92 & 1.27 & 1.52 to 2.16 \\
11:00 & 865 & 0.46 & 0.56 & 4.20 & 3.58 & 1.33 & 0.73 & 1.13 to 1.75 \\
11:30 & 865 & 2.19 & 1.61 & 3.95 & 3.33 & 1.97 & 1.36 & 1.48 to 1.97 \\
12:00 & 866 & 1.12 & 1.13 & 3.17 & 2.58 & 1.20 & 0.62 & 0.82 to 1.44 \\
12:30 & 866 & 1.78 & 0.64 & 3.30 & 2.72 & 0.91 & 0.34 & 0.71 to 0.91 \\
13:00 & 866 & 0.23 & 0.96 & 2.66 & 2.09 & 1.37 & 0.81 & 0.80 to 1.37 \\
13:30 & 866 & 0.79 & 1.02 & 2.51 & 1.97 & 1.63 & 1.09 & 1.22 to 1.63 \\
14:00 & 866 & 0.50 & 0.49 & 2.04 & 1.36 & 1.04 & 0.38 & 0.50 to 1.16 \\
14:30 & 866 & 1.00 & 0.57 & 2.00 & 1.50 & 0.67 & 0.17 & 0.47 to 0.98 \\
15:00 & 866 & 2.51 & 1.36 & 1.93 & 1.40 & 1.33 & 0.80 & 0.60 to 1.44 \\
15:30 & 866 & 0.20 & 1.34 & 0.24 & \ensuremath{-}0.23 & 1.30 & 0.83 & 0.92 to 1.51 \\
\bottomrule
\end{tabular}
}
\end{table}

% Source: results/atm_straddle_intraday_holdclose/mean_by_entry_hhmm_as.png
%   (holdclose notebook, section 8; censored-implied bars flat).
\begin{figure}[H]
\centering
\treefigorpending{0.9\textwidth}{../results/atm_straddle_intraday_holdclose/mean_by_entry_hhmm_as.png}
\caption{Delta-hedged hold to the close by entry clock, block-diagonal ridge:
mean return (top) and annualized Sharpe ratio (bottom) of always short and
$\mathrm{sign}(s)$ (holdclose notebook, where the bars with a censored implied
volatility sit flat).}
\label{fig:app_run_perclock}
\end{figure}

\paragraph{Finding.} Hedged and held to the close, always short is strongest
in the late morning (\runappPcAsPeak{} at \runappPcAsPeakClock{},
\runappPcAsSecond{} at \runappPcAsSecondClock{}) and falls through the
afternoon to \runappPcAsFifteen{} at 15:00 and \runappPcAsLast{} at 15:30
(\runappPcAsLastX{} crossed). At \runappPcAsBeatsAll{} of the
\runappPcDayClocks{} clocks from 09:35 to 15:00 it is ahead of the
$\mathrm{sign}(s)$ rule of every one of the eight forecasts at the midpoint;
at the crossed spread it is ahead of the ridge's at \runappPcAsBeatsAllX{} of
them. The order reverses only at 15:30: $\mathrm{sign}(s)$ \runappPcSgLast{}
(\runappPcSgLastLo{} to \runappPcSgLastHi{} across the eight), \runappPcSgLastX{}
crossed. At the 09:35 open always short earns \runappPcAsFirst{}
(\runappPcAsFirstX{} crossed) and the ridge's $\mathrm{sign}(s)$
\runappPcSgFirst{} (\runappPcSgFirstX{}). In the unhedged re-pick the
always-short Sharpe ratio moves between \runappPcRpAsLo{} and \runappPcRpAsHi{}
from clock to clock, and $\mathrm{sign}(s)$ is ahead at \runappPcRpSgBeats{}
of \runappPcRpClocks{} clocks, all before the spread.

\paragraph{Why it is not in the main text.} What varies through the day is the
premium earned by the forecast-free control, a different object from the
one-bar forecast the paper studies; the forecast is ahead of that control only
at the last bar, which is the main text's trade.

% ----------------------------------------------------------------------------
\subsection{A Forecast Built for the Rest of the Day}
\label{sec:app_run_restofday}

\paragraph{Question.} The rest-of-day signals above reach the close from a
one-bar forecast. Does a model fitted directly on the rest of the day forecast
it better? For every entry clock $c$ from 09:30 to 15:30 the per-bar linear
estimators (ridge, lasso, elastic net; 2000-session rolling window, causal
penalty tuning) are fitted with the realized variance from $c$ to the 16:00
close as the target, on the three column sets of the per-bar models. Each is
compared on identical rows with the one-step construction of the same
estimator and column set: the forecast of the bar starting at $c$ divided by
that bar's prior-days share of the rest of the day. The direct forecast is
back-transformed plainly --- the squared prediction times the arm's causal
per-clock profile of the rest-of-day variance --- with no second-moment
correction.

% Source: results/linear_subsection_restofday/score_rest_of_day.csv, the "plain" columns
%   (qlike_plain, pct_plain_vs_onestep, p_plain_vs_onestep); experiments/score_rest_of_day.py,
%   campaign specs/causal_tune_rest_of_day.py.  The scorer's "causal" column adds a
%   correction to the direct side only and is not used.
\begin{table}[H]
\centering
\caption{Direct rest-of-day forecasts against the one-step construction,
\runappRdNlo{}--\runappRdNhi{} sessions from \runappRdFirst{} to
\runappRdLast{}. $\Delta$QLIKE is $100\,(\mathrm{QLIKE}_{\text{direct}}/
\mathrm{QLIKE}_{\text{one-step}}-1)$ at one entry clock; the median is over
the twelve clocks 09:30--15:00. Clocks better (worse) count the clocks at
which the direct forecast's loss is lower (higher) with a Diebold--Mariano
$p$ below \runappRdAlpha{}; clocks higher counts point estimates above zero.
At 15:30 the two constructions are the same forecast.}
\label{tab:app_run_restofday}
\runappfit{\small% AUTO-GENERATED by writeup/make_appendix_running_tex.py -- do not edit.
% Source: results/linear_subsection_restofday/score_rest_of_day.csv ("plain" columns: qlike_plain, pct_plain_vs_onestep, p_plain_vs_onestep)
\begin{tabular}{llrrrrrr}
\toprule
column set & estimator & sessions & \shortstack{median\\$\Delta$QLIKE (\%)} & \shortstack{clocks\\better} & \shortstack{clocks\\worse} & \shortstack{clocks higher\\(point)} & \shortstack{15:30\\$|\Delta|$ (\%)} \\
\midrule
HAR + calendar & ridge & 1{,}404--1{,}406 & \ensuremath{-}18.8 & 11 & 0 & 0 & 0.000 \\
HAR + calendar & lasso & 1{,}404--1{,}406 & \ensuremath{-}18.0 & 11 & 0 & 0 & 0.000 \\
HAR + calendar & elastic net & 1{,}404--1{,}406 & \ensuremath{-}17.8 & 11 & 0 & 0 & 0.000 \\
\midrule
live-feasible (16 columns) & ridge & 1{,}404--1{,}406 & \ensuremath{-}30.6 & 6 & 0 & 2 & 0.000 \\
live-feasible (16 columns) & lasso & 1{,}404--1{,}406 & \ensuremath{-}31.2 & 9 & 0 & 1 & 0.000 \\
live-feasible (16 columns) & elastic net & 1{,}404--1{,}406 & \ensuremath{-}36.7 & 10 & 0 & 1 & 0.000 \\
\midrule
all features & ridge & 1{,}404--1{,}406 & \ensuremath{-}39.6 & 8 & 0 & 1 & 0.000 \\
all features & lasso & 1{,}404--1{,}406 & \ensuremath{-}36.3 & 9 & 0 & 1 & 0.002 \\
all features & elastic net & 1{,}404--1{,}406 & \ensuremath{-}44.4 & 8 & 0 & 1 & 0.018 \\
\bottomrule
\end{tabular}
}
\end{table}

\paragraph{Finding.} The direct forecast has the lower loss at the median
clock for every estimator and column set: \runappRdBaseLo{}\% to
\runappRdBaseHi{}\% on HAR + calendar, \runappRdLiveLo{}\% to
\runappRdLiveHi{}\% on the live-feasible set and \runappRdAllLo{}\% to
\runappRdAllHi{}\% on all features. It is significantly better at
\runappRdBetterLo{} to \runappRdBetterHi{} of the twelve clocks and never
significantly worse; at 15:30 the two agree to \runappRdLastMax{}\%. The plain
back-transform has rare extreme values: \runappRdExtremeN{} of the
\runappRdCells{} clock cells carry a QLIKE above 1 on one side or the other,
the largest the \runappRdExtremeCell{} at \runappRdExtremeClock{}
(\runappRdExtremeQ{} against \runappRdExtremeOne{} for the one-step
construction), and the same kind of value gives the all-features ridge a plain
QLIKE of \runappRdAllRidgeLast{} at 15:30. None of those cells is a significant
comparison, which is why the table reports the median clock and per-clock tests
rather than a pooled mean.

% The trading test of the direct forecast (checklist item E2) belongs to another workstream:
% experiments/restofday_trading_test.py -> results/linear_subsection_restofday/trading_test/
% (SUMMARY.md, trading_test_*.csv) and its table generated/appendix_running_restofday_trading.tex.
% It landed at 02:34 on 2026-09-29, while this appendix was being assembled; the prose numbers
% below are read from its trading_test_pooled.csv by make_appendix_running_tex.py into
% generated/appendix_running_restofday_trade_numbers.tex.  Both hooks vanish if the outputs do.
\IfFileExists{generated/appendix_running_restofday_trading.tex}{%
\IfFileExists{generated/appendix_running_restofday_trade_numbers.tex}{%
% AUTO-GENERATED by writeup/make_appendix_running_tex.py -- do not edit.
% Source: results/linear_subsection_restofday/trading_test/trading_test_pooled.csv (pool 10:00-15:30; per-bar ridge, plain back-transform unless counted over all)
\newcommand{\runappTtLfCur}{2.12}
\newcommand{\runappTtLfDir}{2.49}
\newcommand{\runappTtLfD}{\ensuremath{+}0.37}
\newcommand{\runappTtLfCi}{[\ensuremath{-}0.38, 1.05]}
\newcommand{\runappTtLfCurX}{1.00}
\newcommand{\runappTtLfDirX}{1.48}
\newcommand{\runappTtLfDX}{\ensuremath{+}0.48}
\newcommand{\runappTtLfCiX}{[\ensuremath{-}0.20, 1.20]}
\newcommand{\runappTtLfBuyCur}{43.6}
\newcommand{\runappTtLfBuyDir}{37.6}
\newcommand{\runappTtLfQ}{\ensuremath{-}26.6}
\newcommand{\runappTtLfVsShort}{\ensuremath{-}0.44}
\newcommand{\runappTtLfVsShortCi}{[\ensuremath{-}2.02, 0.86]}
\newcommand{\runappTtAfCur}{2.23}
\newcommand{\runappTtAfDir}{2.54}
\newcommand{\runappTtAfD}{\ensuremath{+}0.31}
\newcommand{\runappTtAfCi}{[\ensuremath{-}0.37, 0.97]}
\newcommand{\runappTtAfCurX}{1.13}
\newcommand{\runappTtAfDirX}{1.58}
\newcommand{\runappTtAfDX}{\ensuremath{+}0.45}
\newcommand{\runappTtAfCiX}{[\ensuremath{-}0.19, 1.16]}
\newcommand{\runappTtAfBuyCur}{44.2}
\newcommand{\runappTtAfBuyDir}{37.9}
\newcommand{\runappTtAfQ}{\ensuremath{-}62.1}
\newcommand{\runappTtAfVsShort}{\ensuremath{-}0.40}
\newcommand{\runappTtAfVsShortCi}{[\ensuremath{-}1.99, 0.96]}
\newcommand{\runappTtBaCur}{1.33}
\newcommand{\runappTtBaDir}{1.45}
\newcommand{\runappTtBaD}{\ensuremath{+}0.12}
\newcommand{\runappTtBaCi}{[\ensuremath{-}0.63, 0.85]}
\newcommand{\runappTtBaCurX}{0.20}
\newcommand{\runappTtBaDirX}{0.40}
\newcommand{\runappTtBaDX}{\ensuremath{+}0.20}
\newcommand{\runappTtBaCiX}{[\ensuremath{-}0.51, 0.96]}
\newcommand{\runappTtBaBuyCur}{46.2}
\newcommand{\runappTtBaBuyDir}{44.6}
\newcommand{\runappTtBaQ}{\ensuremath{-}18.1}
\newcommand{\runappTtBaVsShort}{\ensuremath{-}1.49}
\newcommand{\runappTtBaVsShortCi}{[\ensuremath{-}3.25, 0.03]}
\newcommand{\runappTtRidgeIntervals}{6}
\newcommand{\runappTtRidgeZero}{6}
\newcommand{\runappTtCombos}{18}
\newcommand{\runappTtAhead}{15}
\newcommand{\runappTtExclFor}{0}
\newcommand{\runappTtExclAgainst}{0}
%
\paragraph{Trading test.} The direct forecast replaces the stretch in the
delta-hedged hold-to-close signal of Section~\ref{sec:app_run_holdclose},
$s_t=F^{\mathrm{rem}}_t-\mathrm{IV}^2_{\mathrm{hr},t}h_t$, on the same days,
bars, fills and hedge (Table~\ref{tab:app_run_restofday_trade}; per-bar ridge,
plain back-transform). Pooled over the entry clocks 10:00--15:30 it trades
better by a point estimate: on the live-feasible set from \runappTtLfCur{} to
\runappTtLfDir{} at the midpoint (\runappTtLfD{} \runappTtLfCi{}) and from
\runappTtLfCurX{} to \runappTtLfDirX{} crossed (\runappTtLfDX{}
\runappTtLfCiX{}); on all features \runappTtAfD{} \runappTtAfCi{}; on HAR +
calendar \runappTtBaD{} \runappTtBaCi{}. Of these \runappTtRidgeIntervals{}
intervals (three column sets, two fills), \runappTtRidgeZero{} include zero.
Over all \runappTtCombos{}
estimator, column-set and back-transform combinations at the midpoint the
direct forecast is ahead in \runappTtAhead{}, with \runappTtExclFor{} intervals
excluding zero in its favour and \runappTtExclAgainst{} against. It buys less
often (\runappTtLfBuyDir{}\% of bars against \runappTtLfBuyCur{}\% on the
live-feasible set), the mirror image of the fitted maps of
Table~\ref{tab:app_run_maps}, although it changes the QLIKE on these bars by
\runappTtLfQ{}\%; both signals stay below always short (direct minus always
short \runappTtLfVsShort{} \runappTtLfVsShortCi{}, live-feasible, midpoint).
The trading gain, where there is one, is inside its sampling error.
}{}%
\begin{table}[H]
\centering
\caption{Trading test of the direct rest-of-day forecast in the delta-hedged
hold-to-close book, per-bar ridge, on the \runappHcNdays{} days (another
workstream's table, read as generated). Top: pooled over the entry clocks
10:00--15:30 (daily sums), Sharpe ratios of the current signal
($\widehat{RV}_t/w_t$) and the direct one ($F^{\mathrm{rem}}_t$) and their
paired difference, for the plain back-transform and a recalibration per entry
clock. Bottom: by entry clock, plain, midpoint. Intervals: paired circular
block bootstrap, 95\% percentile; $^{*}$: the percentile and basic intervals
both exclude zero.}
\label{tab:app_run_restofday_trade}
% The generated file holds two tabulars separated by vertical space, and its first is
% wider than the text (2026-09-29 build: by 62.8pt).  A minipage 15% wider than the text
% holds both; the whole is then scaled back to the text width.
\resizebox{\textwidth}{!}{\begin{minipage}{1.15\textwidth}\centering
% AUTO-GENERATED by experiments/restofday_trading_test.py -- do not edit.
% Source: results/linear_subsection_restofday/trading_test/trading_test_{pooled,by_entry_time,reference_rules}.csv
% Table 1: pooled 10:00--15:30 (daily sums), per-bar ridge. Table 2: by entry time, plain back-transform, midpoint.
\begingroup\small\setlength{\tabcolsep}{3pt}\noindent
\begin{tabular}{llrrlrrlrr}
\toprule
& & \multicolumn{3}{c}{midpoint} & \multicolumn{3}{c}{crossed spread} & \multicolumn{2}{c}{buy \%} \\
\cmidrule(lr){3-5}\cmidrule(lr){6-8}\cmidrule(lr){9-10}
bucket & $F$ & current & direct & direct $-$ current & current & direct & direct $-$ current & cur. & dir. \\
\midrule
\multicolumn{10}{l}{\emph{No rest-of-day forecast: always short 2.93 / 2.21; block-diagonal ridge $\mathrm{sign}(s)$ 2.40 / 1.35 (midpoint / crossed)}} \\
\midrule
\multicolumn{10}{l}{\emph{plain back-transform}} \\
live-feasible & plain & 2.12 & 2.49 & +0.37 [-0.38, +1.05] & 1.00 & 1.48 & +0.48 [-0.20, +1.20] & 44 & 38 \\
all features & plain & 2.23 & 2.54 & +0.31 [-0.37, +0.97] & 1.13 & 1.58 & +0.45 [-0.19, +1.16] & 44 & 38 \\
HAR + calendar & plain & 1.33 & 1.45 & +0.12 [-0.63, +0.85] & 0.20 & 0.40 & +0.20 [-0.51, +0.96] & 46 & 45 \\
\multicolumn{10}{l}{\emph{recalibrated per entry time}} \\
live-feasible & recal & 1.40 & 1.82 & +0.42 [-0.30, +1.18] & 0.25 & 0.80 & +0.55 [-0.14, +1.39] & 52 & 46 \\
all features & recal & 1.33 & 1.81 & +0.48 [-0.25, +1.28] & 0.21 & 0.86 & +0.65 [-0.06, +1.53] & 52 & 46 \\
HAR + calendar & recal & 0.41 & 0.21 & -0.20 [-1.00, +0.60] & -0.73 & -0.83 & -0.10 [-0.85, +0.72] & 54 & 55 \\
\bottomrule
\end{tabular}
\endgroup

\medskip

\begingroup\scriptsize\setlength{\tabcolsep}{3pt}\noindent
\resizebox{\textwidth}{!}{%
\begin{tabular}{lrrrlrrlrrl}
\toprule
& always & \multicolumn{3}{c}{live-feasible} & \multicolumn{3}{c}{all features} & \multicolumn{3}{c}{HAR + calendar} \\
\cmidrule(lr){3-5}\cmidrule(lr){6-8}\cmidrule(lr){9-11}
entry & short & current & direct & direct $-$ current & current & direct & direct $-$ current & current & direct & direct $-$ current \\
\midrule
09:35 & 3.11 & 1.23 & 1.04 & -0.19 [-1.29, +0.86] & 1.35 & 1.06 & -0.29 [-1.35, +0.84] & 1.04 & 1.48 & +0.44 [-0.77, +1.54] \\
10:00 & 3.32 & 0.53 & 1.79 & +1.26 [+0.13, +2.51]$^{*}$ & 0.64 & 1.60 & +0.96 [-0.13, +2.00] & -0.07 & 0.85 & +0.92 [-0.07, +2.02] \\
10:30 & 4.18 & 0.46 & 0.88 & +0.42 [-0.64, +1.55] & 0.53 & 1.49 & +0.96 [-0.14, +2.02] & 0.15 & 0.91 & +0.76 [-0.31, +1.82] \\
11:00 & 4.20 & 0.85 & 1.84 & +0.99 [-0.17, +2.20] & 0.83 & 1.78 & +0.95 [-0.13, +2.08] & 0.76 & 0.93 & +0.16 [-1.05, +1.32] \\
11:30 & 3.95 & 0.84 & 1.13 & +0.28 [-0.89, +1.51] & 0.11 & 1.52 & +1.41 [+0.29, +2.68]$^{*}$ & 0.25 & 0.35 & +0.10 [-1.07, +1.26] \\
12:00 & 3.17 & 0.84 & 0.37 & -0.47 [-1.60, +0.62] & 0.85 & 1.04 & +0.18 [-1.02, +1.40] & 0.50 & 0.26 & -0.24 [-1.31, +0.87] \\
12:30 & 3.30 & 1.20 & 1.29 & +0.09 [-0.83, +1.10] & 1.54 & 1.16 & -0.39 [-1.33, +0.62] & 0.84 & 0.76 & -0.08 [-1.18, +1.00] \\
13:00 & 2.66 & 1.24 & 1.42 & +0.18 [-0.94, +1.32] & 1.48 & 1.19 & -0.30 [-1.31, +0.78] & 1.47 & 0.42 & -1.04 [-2.01, -0.10]$^{*}$ \\
13:30 & 2.51 & 1.27 & 1.52 & +0.24 [-0.89, +1.42] & 1.42 & 1.76 & +0.35 [-0.66, +1.39] & 0.71 & 0.71 & -0.00 [-1.27, +1.15] \\
14:00 & 2.04 & 2.09 & 1.85 & -0.24 [-1.24, +0.76] & 1.74 & 1.95 & +0.21 [-0.78, +1.23] & 1.15 & 1.10 & -0.05 [-0.96, +0.84] \\
14:30 & 2.00 & 0.68 & 1.15 & +0.47 [-0.15, +1.15] & 1.06 & 1.22 & +0.15 [-0.54, +0.85] & 0.34 & 0.90 & +0.57 [-0.29, +1.43] \\
15:00 & 1.93 & 0.35 & 1.43 & +1.08 [+0.38, +1.85]$^{*}$ & 0.74 & 1.53 & +0.79 [-0.01, +1.60] & -0.00 & 0.45 & +0.45 [-0.26, +1.22] \\
15:30 & 0.24 & 1.79 & 1.79 & same forecast & 1.85 & 1.85 & same forecast & 1.31 & 1.31 & same forecast \\
\bottomrule
\end{tabular}}
\endgroup

% Per-bar ridge, plain back-transform, midpoint; each row is one entry time's daily series (863--866 days); intervals: paired circular block bootstrap, 95 % percentile; * = the percentile and basic intervals both exclude zero.

\end{minipage}}
\end{table}}{}

\paragraph{Why it is not in the main text.} The paper's target is the one-bar
variance and its rest-of-session sections are parked. A better rest-of-day
loss has also not been shown to be a better trade: the fitted maps of
Table~\ref{tab:app_run_maps} forecast about as well as the stretch and trade
worse.

% ----------------------------------------------------------------------------
\subsection{After the Research Sample: May 2024 to September 2026}
\label{sec:app_run_post2024}

\paragraph{Question.} The main text's scored frame ends with the forecast
panel on 30~April~2024. What did the trade do afterwards? It has been run
forward in the long-only form used live (``the card''): at 15:30 buy the
straddle at the ask when the ask is at or below $P^{*}$, the Black--76 value of
the straddle at the forecast's total volatility $\sqrt{\widehat{RV}_t}$, and
hold it to the close; month-ends are bought at any price and are left to
Section~\ref{sec:app_run_monthend}. With the implied variance read at the ask,
this is the buy side of $\mathrm{sign}(s)$. The forecast is the per-bar ridge
on the free-feed column set of Section~\ref{sec:app_run_vix}, refitted on the
panel extended with the post-2024 feeds (index-futures minute bars from a second
vendor, hourly VIX prints). P\&L per day counts $R$ at the ask on the days
bought and zero on the others. Studies 85 to 90 ask whether the loss after
2024 comes from the inputs (study 87), the buy threshold (88), the training
window or the target (89), or the regime (90).

% Sources: results/atm_straddle_intraday_holdclose/proposals/{85,87,88,89,90}/ (summary.txt,
%   summary.csv, forecast_diagnostics.csv); scripts writeup/intraday_proposals/85_earliest_limit.py,
%   87_daily_edge_inputs.py, 88_walkforward_threshold.py, 89_recent_and_ratio_models.py,
%   90_regime_filter.py.  Study 88 at its N = 250 window; study 89 at the official close.
\begin{table}[H]
\centering
\caption{The card after the research sample, month-ends excluded. Panel (a):
research years 2021-09 to 2024-04 (\runappPoResDays{} days; the first 250
eligible days from 2020-01 are the walk-forward warm-up), after 2024-05
\runappPoPostDays{} days; each walk-forward variant is set from the 250
eligible days before each decision. Panel (b): research years 2020-01 to 2024-04
(\runappPoEightyNineResDays{} days; the target models start later, once their
windows fill), after 2024-05 the same \runappPoPostDays{} days. Panel (c): the
research years only (\runappPoEightySevenDays{} days) with the post-2024 input
changes applied. The last column is the paired difference from the card on the
same days --- after 2024-05 in panels (a) and (b), on the research years in
panel (c) --- with a 95\% circular block-bootstrap interval.}
\label{tab:app_run_post2024}
\runappfit{\small% AUTO-GENERATED by writeup/make_appendix_running_tex.py -- do not edit.
% Source: results/atm_straddle_intraday_holdclose/proposals/88/summary.csv (N = 250)
% Source: results/atm_straddle_intraday_holdclose/proposals/90/summary.csv
% Source: results/atm_straddle_intraday_holdclose/proposals/89/summary.csv (official close)
% Source: results/atm_straddle_intraday_holdclose/proposals/87/summary.csv
\begin{tabular}{lrrrr}
\toprule
 & \shortstack{\% days bought\\research / after} & \shortstack{research\\P\&L/day ($t$)} & \shortstack{after 2024-05\\P\&L/day ($t$)} & \shortstack{minus the card\\{}[95\%]} \\
\midrule
\multicolumn{5}{l}{\emph{(a) the card and walk-forward variants (studies 88 and 90)}} \\
card: buy iff ask $\le P^{*}$ & 38 / 51 & 0.0538 (1.60) & \ensuremath{-}0.0114 (\ensuremath{-}0.35) & --- \\
88: trailing-optimal multiplier on $P^{*}$ & 36 / 36 & 0.0208 (0.77) & 0.0225 (0.77) & \ensuremath{+}0.0339 [0.009, 0.061] \\
88: regression of $R$ on $\log(P^{*}/\text{ask})$ & 14 / 24 & 0.0050 (0.32) & \ensuremath{-}0.0061 (\ensuremath{-}0.27) & \ensuremath{+}0.0052 [\ensuremath{-}0.046, 0.053] \\
88: premium tightened by trailing realized/implied & 32 / 28 & 0.0572 (1.77) & 0.0154 (0.66) & \ensuremath{+}0.0267 [\ensuremath{-}0.027, 0.075] \\
90: VIX tercile & 32 / 26 & 0.0477 (1.53) & \ensuremath{-}0.0092 (\ensuremath{-}0.37) & \ensuremath{+}0.0021 [\ensuremath{-}0.047, 0.041] \\
90: realized variance so far, tercile & 25 / 24 & 0.0126 (0.55) & \ensuremath{-}0.0132 (\ensuremath{-}0.59) & \ensuremath{-}0.0019 [\ensuremath{-}0.053, 0.042] \\
90: FOMC statement day & 30 / 35 & 0.0434 (1.45) & \ensuremath{-}0.0381 (\ensuremath{-}1.45) & \ensuremath{-}0.0268 [\ensuremath{-}0.065, 0.002] \\
90: implied over trailing realized, tercile & 24 / 27 & 0.0523 (1.83) & \ensuremath{-}0.0128 (\ensuremath{-}0.53) & \ensuremath{-}0.0015 [\ensuremath{-}0.041, 0.030] \\
90: the card's last 20 buys & 26 / 23 & 0.0373 (1.27) & 0.0073 (0.31) & \ensuremath{+}0.0187 [\ensuremath{-}0.021, 0.058] \\
\midrule
\multicolumn{5}{l}{\emph{(b) training windows and trade-aligned targets (study 89)}} \\
training window 1000 sessions (card: 2000) & 33 / 57 & 0.0518 (1.92) & \ensuremath{-}0.0190 (\ensuremath{-}0.54) & \ensuremath{-}0.0076 [\ensuremath{-}0.036, 0.020] \\
training window 500 sessions & 38 / 58 & 0.0173 (0.64) & \ensuremath{-}0.0240 (\ensuremath{-}0.67) & \ensuremath{-}0.0126 [\ensuremath{-}0.049, 0.025] \\
training window 250 sessions & 40 / 58 & 0.0223 (0.77) & \ensuremath{-}0.0124 (\ensuremath{-}0.35) & \ensuremath{-}0.0010 [\ensuremath{-}0.035, 0.034] \\
target $\log(RV/\mathrm{IV}^2_{\text{ask}})$, ridge, 250 days & 3 / 2 & 0.0058 (0.54) & 0.0015 (0.22) & \ensuremath{+}0.0129 [\ensuremath{-}0.050, 0.071] \\
target $R>0$, logistic, 250 days & 18 / 18 & \ensuremath{-}0.0010 (\ensuremath{-}0.06) & 0.0015 (0.09) & \ensuremath{+}0.0129 [\ensuremath{-}0.057, 0.077] \\
target $\log(RV/\mathrm{IV}^2_{\text{ask}})$, ridge, 500 days & 5 / 5 & 0.0139 (0.73) & 0.0017 (0.18) & \ensuremath{+}0.0131 [\ensuremath{-}0.050, 0.072] \\
\midrule
\multicolumn{5}{l}{\emph{(c) the research years rerun with the post-2024 inputs (study 87)}} \\
research card (per-bar ridge, live-feasible columns) & 34 & 0.0455 (1.72) & --- & --- \\
free-feed column set & 34 & 0.0594 (2.16) & --- & \ensuremath{+}0.0139 [\ensuremath{-}0.002, 0.035] \\
VIX family as hourly prints & 34 & 0.0472 (1.71) & --- & \ensuremath{+}0.0017 [\ensuremath{-}0.014, 0.022] \\
ES moments rescaled to the new vendor & 35 & 0.0417 (1.58) & --- & \ensuremath{-}0.0038 [\ensuremath{-}0.022, 0.012] \\
all three together & 35 & 0.0360 (1.47) & --- & \ensuremath{-}0.0095 [\ensuremath{-}0.045, 0.022] \\
\bottomrule
\end{tabular}
}
\end{table}

\paragraph{Finding.} The card has no edge after the research sample. On the
SPX book from May 2024 to December 2025 it makes \runappPoSpxPnl{} per day
($t=\runappPoSpxT{}$, \runappPoSpxBuys{} buys on \runappPoSpxDays{} days). In
the walk-forward frame it makes \runappPoResPnl{} per day
($t=\runappPoResT{}$) in the research years, buying on \runappPoResShare{}\%
of days, and \runappPoPostPnl{} ($t=\runappPoPostT{}$, interval
\runappPoPostCi{}) after April 2024, buying on \runappPoPostShare{}\%. The
post-2024 inputs do not explain the loss: applied together to the research
years they leave \runappPoAllPnl{} per day against \runappPoDeckPnl{}
(difference \runappPoAllD{}, interval \runappPoAllCi{}), and each alone moves
the P\&L per day by at most \runappPoAloneMax{}. No walk-forward repair
restores a material edge without giving up the research years: the
trailing-optimal multiplier on $P^{*}$ makes \runappPoFOnePost{} per day after
2024 ($t=\runappPoFOneT{}$, interval \runappPoFOneCi{}), above the card by
\runappPoFOneD{} \runappPoFOneDCi{} (Holm $p=\runappPoFOneHolm{}$), but costs
\runappPoFOneResD{} \runappPoFOneResDCi{} per day in the research years;
every regime filter's post-2024 interval includes zero, the best being the
card's own recent form (\runappPoCFourPost{}, difference interval
\runappPoCFourCi{}); shorter training windows lose \runappPoTwHi{} to
\runappPoTwLo{} per day after 2024; and models of the trade-aligned targets
buy on \runappPoRatioShareLo{}\% to \runappPoRatioShareHi{}\% of days and make
\runappPoRatioPnlLo{} to \runappPoRatioPnlHi{}. After 2024 the variance implied
at the ask forecasts the 15:30--16:00 realized variance better than the model
does (QLIKE \runappPoQAsk{} against \runappPoQModelLo{} to
\runappPoQModelHi{} across the training windows); in the research years the
order was the reverse (\runappPoQCardRes{} against \runappPoQAskRes{}).

\paragraph{Why it is not in the main text.} It tests the live, long-only form
of the trade on a rebuilt panel and on days the forecast panel of the main
text does not cover. It is the trade's only record after April 2024.

% ----------------------------------------------------------------------------
\subsection{The Month-End Long Straddle}
\label{sec:app_run_monthend}

\paragraph{Question.} On the last session of each month the card buys the
straddle at any price, with no forecast. Is 15:30 the right clock for that
purchase? Study 83 buys, at each clock, the straddle nearest the spot at the
ask and settles it at the close.

% Source: results/atm_straddle_intraday_holdclose/proposals/83/summary.txt, gated in the maker
%   against 83/month_end_by_clock.csv (mean R per clock); script
%   writeup/intraday_proposals/83_month_end_entry_clock.py.
\begin{table}[H]
\centering
\caption{The month-end long straddle by entry clock, \runappMeN{} month-ends
from \runappMeFirst{} to \runappMeLast{}, bought at the ask. The spread is the
median relative spread at entry; minus 15:30 is the paired difference in mean
$R$, with its paired $t$ and the Holm-adjusted $p$ over the twelve earlier
clocks.}
\label{tab:app_run_monthend}
\runappfit{\small% AUTO-GENERATED by writeup/make_appendix_running_tex.py -- do not edit.
% Source: results/atm_straddle_intraday_holdclose/proposals/83/summary.txt
% Source: results/atm_straddle_intraday_holdclose/proposals/83/month_end_by_clock.csv (gate: mean R)
\begin{tabular}{lrrrrrrrr}
\toprule
entry & \shortstack{median\\ask} & \shortstack{median\\spread (\%)} & mean $R$ & $t$ & hit (\%) & minus 15:30 & paired $t$ & Holm $p$ \\
\midrule
09:35 & 23.80 & 1.9 & 0.152 & 1.41 & 46 & \ensuremath{-}0.281 & \ensuremath{-}1.54 & 1.00 \\
10:00 & 22.75 & 3.2 & 0.211 & 1.84 & 47 & \ensuremath{-}0.222 & \ensuremath{-}1.25 & 1.00 \\
10:30 & 20.10 & 2.2 & 0.190 & 1.61 & 54 & \ensuremath{-}0.242 & \ensuremath{-}1.39 & 1.00 \\
11:00 & 18.80 & 2.2 & 0.190 & 1.58 & 53 & \ensuremath{-}0.242 & \ensuremath{-}1.51 & 1.00 \\
11:30 & 18.30 & 2.3 & 0.202 & 1.76 & 51 & \ensuremath{-}0.230 & \ensuremath{-}1.46 & 1.00 \\
12:00 & 16.40 & 2.4 & 0.238 & 1.98 & 54 & \ensuremath{-}0.194 & \ensuremath{-}1.28 & 1.00 \\
12:30 & 15.05 & 2.5 & 0.309 & 2.27 & 54 & \ensuremath{-}0.124 & \ensuremath{-}0.75 & 1.00 \\
13:00 & 15.15 & 2.6 & 0.269 & 1.98 & 51 & \ensuremath{-}0.163 & \ensuremath{-}1.04 & 1.00 \\
13:30 & 14.25 & 2.9 & 0.301 & 2.35 & 51 & \ensuremath{-}0.132 & \ensuremath{-}0.96 & 1.00 \\
14:00 & 13.15 & 2.9 & 0.422 & 2.99 & 53 & \ensuremath{-}0.011 & \ensuremath{-}0.08 & 1.00 \\
14:30 & 11.70 & 2.9 & 0.355 & 2.50 & 56 & \ensuremath{-}0.077 & \ensuremath{-}0.67 & 1.00 \\
15:00 & 9.95 & 4.0 & 0.267 & 1.82 & 53 & \ensuremath{-}0.166 & \ensuremath{-}1.56 & 1.00 \\
15:30 & 8.30 & 4.9 & 0.432 & 2.62 & 54 & --- & --- & --- \\
\bottomrule
\end{tabular}
}
\end{table}

\paragraph{Finding.} The purchase pays at every clock and most at 15:30 (mean
$R$ \runappMeR{}, $t=\runappMeT{}$); the earlier clocks earn \runappMeEarlyLo{}
to \runappMeEarlyHi{}, and no difference from 15:30 survives the Holm
adjustment (smallest $p$ \runappMeHolmMin{}). The 15:30 return is
\runappMeFirstHalf{} over the first half of the month-ends and
\runappMeSecondHalf{} over the second, from \runappMeSplit{}. On the
\runappDsMonthEnds{} SPX month-ends from December 2024 of
Table~\ref{tab:app_run_decision86} the 15:30 purchase earns
\runappDsMeLast{} at the ask ($t=\runappDsMeLastT{}$), and entries from 15:40
to 15:50 are \runappDsMeLaterLo{} to \runappDsMeLaterHi{} higher, with Holm
$p$ of \runappDsMeHolmMin{}.

\paragraph{Why it is not in the main text.} It is a calendar trade that reads
no forecast.

% ----------------------------------------------------------------------------
\subsection{Is 15:30 the Right Decision Time?}
\label{sec:app_run_decision}

\paragraph{Question.} Two studies move the decision. Study 80 reruns the
per-bar ridge with every regressor lagged $h$ bars, so that the forecast of
the 15:30--16:00 bar uses information through 15:30 ($h=1$, the forecast the
card uses), 15:00 ($h=2$) and so on, and trades the card rule at the 15:30 ask
on the \runappDtDays{} days. Study 86 moves decision and entry together to
$\tau$ between 15:00 and 15:50 on the SPX book of \runappDsDays{} sessions
from \runappDsFirst{} to \runappDsLast{}: a session-by-session regression on
the thirty-minute design moved to $\tau$ forecasts the variance from $\tau$ to
the close, and the card buys at $\tau$ at the ask when the ask is at or below
$P^{*}$.

% Source: results/atm_straddle_intraday_holdclose/proposals/80/summary_by_horizon.csv and
%   bootstrap_diff.csv (scope all_days); script writeup/intraday_proposals/80_earliest_decision_clock.py.
\begin{table}[H]
\centering
\caption{The decision clock on the \runappDtDays{} days (study 80): the card
rule at the SPXW 15:30 ask with the forecast built from information through
the last input. QLIKE is that forecast's loss on the 16:00 bar. The last
column is the paired difference in P\&L per day from $h=1$, with a 95\% block
bootstrap interval (blocks of 20 days).}
\label{tab:app_run_decision80}
\runappfit{\small% AUTO-GENERATED by writeup/make_appendix_running_tex.py -- do not edit.
% Source: results/atm_straddle_intraday_holdclose/proposals/80/summary_by_horizon.csv (scope all_days)
% Source: results/atm_straddle_intraday_holdclose/proposals/80/bootstrap_diff.csv (scope all_days)
\begin{tabular}{lrrrrrrr}
\toprule
$h$ & \shortstack{last\\input} & \shortstack{QLIKE\\16:00 bar} & \shortstack{days\\bought} & \shortstack{mean $R$\\bought} & $t$ & P\&L/day & \shortstack{minus $h=1$\\{}[95\%]} \\
\midrule
1 & 15:30 & 0.1054 & 286 & 0.152 & 1.97 & 0.0501 & --- \\
2 & 15:00 & 0.1209 & 341 & 0.010 & 0.17 & 0.0040 & \ensuremath{-}0.046 [\ensuremath{-}0.082, \ensuremath{-}0.016] \\
3 & 14:30 & 0.1264 & 378 & 0.024 & 0.37 & 0.0105 & \ensuremath{-}0.040 [\ensuremath{-}0.074, \ensuremath{-}0.004] \\
4 & 14:00 & 0.1467 & 403 & 0.015 & 0.24 & 0.0069 & \ensuremath{-}0.043 [\ensuremath{-}0.077, \ensuremath{-}0.011] \\
\bottomrule
\end{tabular}
}
\end{table}

% Source: results/atm_straddle_intraday_holdclose/proposals/86/spx_summary_tau.csv (model T0),
%   spx_bootstrap_tau.csv (T0, pnl_ask) and spx_month_end_tau.csv; the session range from
%   86/summary.txt; script writeup/intraday_proposals/86_entry_1500_1520.py.
\begin{table}[H]
\centering
\caption{Decision and entry at $\tau$ on the SPX book (study 86), month-ends
excluded, at the ask. The difference from 15:30 is paired, with a 95\% block
bootstrap interval (blocks of 20 sessions) and the Holm-adjusted one-sided $p$
over the ten other clocks. The month-end columns buy every one of the
\runappDsMonthEnds{} month-ends at $\tau$.}
\label{tab:app_run_decision86}
\runappfit{\small% AUTO-GENERATED by writeup/make_appendix_running_tex.py -- do not edit.
% Source: results/atm_straddle_intraday_holdclose/proposals/86/spx_summary_tau.csv (model T0)
% Source: results/atm_straddle_intraday_holdclose/proposals/86/spx_bootstrap_tau.csv (T0, pnl_ask)
% Source: results/atm_straddle_intraday_holdclose/proposals/86/spx_month_end_tau.csv
\begin{tabular}{lrrrrrrr}
\toprule
$\tau$ & \shortstack{days\\bought} & P\&L/day & $t$ & \shortstack{minus 15:30\\{}[95\%]} & Holm $p$ & \shortstack{month-end\\mean $R$} & \shortstack{month-end\\$t$} \\
\midrule
15:00 & 208 & 0.0223 & 0.60 & \ensuremath{+}0.033 [\ensuremath{-}0.042, 0.111] & 1.00 & 0.320 & 1.52 \\
15:05 & 207 & \ensuremath{-}0.0075 & \ensuremath{-}0.22 & \ensuremath{+}0.003 [\ensuremath{-}0.082, 0.092] & 1.00 & 0.173 & 0.76 \\
15:10 & 208 & \ensuremath{-}0.0251 & \ensuremath{-}0.78 & \ensuremath{-}0.015 [\ensuremath{-}0.078, 0.050] & 1.00 & 0.282 & 1.26 \\
15:15 & 214 & 0.0110 & 0.28 & \ensuremath{+}0.021 [\ensuremath{-}0.045, 0.094] & 1.00 & 0.310 & 1.36 \\
15:20 & 221 & \ensuremath{-}0.0318 & \ensuremath{-}0.87 & \ensuremath{-}0.022 [\ensuremath{-}0.072, 0.034] & 1.00 & 0.249 & 1.06 \\
15:25 & 223 & \ensuremath{-}0.0034 & \ensuremath{-}0.09 & \ensuremath{+}0.007 [\ensuremath{-}0.037, 0.052] & 1.00 & 0.250 & 1.06 \\
15:30 & 227 & \ensuremath{-}0.0103 & \ensuremath{-}0.27 & --- & --- & 0.356 & 1.45 \\
15:35 & 210 & 0.0057 & 0.15 & \ensuremath{+}0.016 [\ensuremath{-}0.054, 0.087] & 1.00 & 0.381 & 1.44 \\
15:40 & 190 & 0.0068 & 0.18 & \ensuremath{+}0.017 [\ensuremath{-}0.053, 0.087] & 1.00 & 0.462 & 1.76 \\
15:45 & 192 & 0.0425 & 1.03 & \ensuremath{+}0.053 [\ensuremath{-}0.032, 0.140] & 1.00 & 0.481 & 2.12 \\
15:50 & 156 & 0.0059 & 0.15 & \ensuremath{+}0.016 [\ensuremath{-}0.054, 0.087] & 1.00 & 0.481 & 1.63 \\
\bottomrule
\end{tabular}
}
\end{table}

\paragraph{Finding.} Only the 15:30 forecast trades. On the \runappDtDays{}
days the card's buys return \runappDtOneR{} ($t=\runappDtOneT{}$,
\runappDtOneN{} buys) with the 15:30 forecast and \runappDtTwoR{}
($t=\runappDtTwoT{}$) with the 15:00 one; every earlier forecast earns less per
day with the interval below zero ($h=2$: \runappDtTwoD{} \runappDtTwoCi{}). On
SPX after 2024 no $\tau$ is distinguishable from 15:30 (smallest Holm $p$
\runappDsHolmMin{} over both forecast designs and both fills) and no $\tau$'s
P\&L per day has $|t|$ above \runappDsTmax{}; the SPX spread is
\runappDsSpreadLo{}\% to \runappDsSpreadHi{}\% of premium across the clocks,
against \runappDsXspSpreadLo{}\% to \runappDsXspSpreadHi{}\% for the
one-tenth-size XSP contract. Two further checks sit here. Adding one-minute
index-futures features to the 15:30 forecast (study 82, \runappDgDays{}
out-of-sample sessions) changes its QLIKE by \runappDgFineLo{}\% to
\runappDgFineHi{}\% with $|t|\le\runappDgFineDm{}$. On \runappLoDays{} days
with minute quotes, a limit order at $P^{*}$ cancelled at once earns
\runappLoOnce{} per day ($t=\runappLoOnceT{}$), and leaving it working one
minute longer costs \runappLoOneMin{} per day ($t=\runappLoOneMinT{}$), because
the fills that come late are the ones that lose (study 84).

\paragraph{Why it is not in the main text.} These confirm the timing the main
text uses and settle execution details of the live form.

% ----------------------------------------------------------------------------
\subsection{The VIX Alone}
\label{sec:app_run_vix}

\paragraph{Question.} A live forecaster must fetch every input each afternoon.
Which columns does the forecast need? Four column sets are fitted the way the
per-bar models are (one regression per bar, 2000-session rolling window;
ridge, lasso and elastic net), each on the HAR lags and calendar columns plus:
the live-feasible set, the sixteen design columns \runappKeyLive{}; that set
without \texttt{vvix} and \texttt{vix3m}; the free feed, the live-feasible
set without \runappKeyFreeDrop{} (thirteen columns, all available from free
sources); and \texttt{vix} alone.

% Source: results/atm_straddle_0dte_1530/vix_bucket.csv (notebooks/atm_straddle_rv_iv.executed.ipynb,
%   section 17), research-scorer columns; the key from src/data/loading.py SUBGROUPS
%   (live_feasible, live_vix_only, free_vix_only, vix_only).
\begin{table}[H]
\centering
\caption{Column sets for the 15:30 forecast, on the \runappVxDays{} days.
QLIKE is relative to the live-feasible set on the same rows, on the 16:00 bar
and averaged over the 13 bars of the session; negative is better. The Sharpe
ratios are those of $\mathrm{sign}(s)$ under the research scorer, which
recalibrates the 16:00 bar on its own; the notebook's all-bar recalibration is
quoted in the text.}
\label{tab:app_run_vix}
\runappfit{\small% AUTO-GENERATED by writeup/make_appendix_running_tex.py -- do not edit.
% Source: results/atm_straddle_0dte_1530/vix_bucket.csv (research-scorer columns only; the notebook-scorer Sharpe ratios are quoted in the text)
\begin{tabular}{llrrrr}
\toprule
column set & estimator & \shortstack{16:00 QLIKE vs\\live-feasible (\%)} & \shortstack{13-bar QLIKE vs\\live-feasible (\%)} & \shortstack{Sharpe\\(mid)} & \shortstack{Sharpe\\(crossed)} \\
\midrule
HAR + calendar + VIX only & ridge & \ensuremath{-}5.92 & \ensuremath{+}3.31 & 1.48 & 1.00 \\
HAR + calendar + VIX only & lasso & \ensuremath{-}1.73 & \ensuremath{+}3.11 & 1.42 & 0.94 \\
HAR + calendar + VIX only & elastic net & \ensuremath{-}0.60 & \ensuremath{+}3.34 & 1.34 & 0.86 \\
\midrule
live-feasible minus VVIX and VIX3M & ridge & \ensuremath{-}0.04 & \ensuremath{-}0.44 & 1.86 & 1.39 \\
live-feasible minus VVIX and VIX3M & lasso & \ensuremath{+}0.10 & \ensuremath{-}0.10 & 1.85 & 1.38 \\
live-feasible minus VVIX and VIX3M & elastic net & \ensuremath{+}0.14 & \ensuremath{-}0.51 & 1.90 & 1.43 \\
\midrule
live-feasible (16 columns) & ridge & 0.00 & 0.00 & 1.90 & 1.44 \\
live-feasible (16 columns) & lasso & 0.00 & 0.00 & 1.83 & 1.36 \\
live-feasible (16 columns) & elastic net & 0.00 & 0.00 & 1.41 & 0.94 \\
\midrule
free feed (ES 1-min bars + VIX + FOMC calendar) & ridge & \ensuremath{-}0.43 & \ensuremath{-}0.01 & 1.83 & 1.36 \\
free feed (ES 1-min bars + VIX + FOMC calendar) & lasso & \ensuremath{-}0.63 & \ensuremath{+}0.76 & 1.91 & 1.45 \\
free feed (ES 1-min bars + VIX + FOMC calendar) & elastic net & \ensuremath{+}0.11 & \ensuremath{+}0.18 & 1.81 & 1.35 \\
\bottomrule
\end{tabular}
}
\end{table}

\paragraph{Finding.} The VIX alone gives the best 16:00 forecast (QLIKE
\runappVxQ{}\% for the ridge) but is \runappVxQThirteen{}\% worse over the 13
bars, and it trades worse: \runappVxMid{} at the midpoint and \runappVxX{}
crossed against \runappVxLfMid{} and \runappVxLfX{}. Under the notebook's
all-bar recalibration the order is the same (\runappVxNbMid{} against
\runappVxNbLfMid{}). Dropping \texttt{vvix} and \texttt{vix3m} changes the
16:00 QLIKE by \runappVxMinusQ{}\%, and the free feed by \runappVxFreeQ{}\%
(Sharpe \runappVxFreeMid{}): the index-futures moments are what the trade uses,
and of the Cboe indices only the VIX is needed.

\paragraph{Why it is not in the main text.} It is an implementation question
for the live forecast. It is also one more case of the loss and the trade
disagreeing.

% ----------------------------------------------------------------------------
\subsection{Per-Bar Against Pooled Forecasts}
\label{sec:app_run_perbar}

\paragraph{Question.} A per-bar forecast fits one regression for each
thirty-minute bar; its pooled twin fits one regression for all bars. Which is
the better variance forecast? QLIKE is computed bar by bar for the 13 session
bars on the \runappPbDays{} days, both tables through the same recalibration.

% Source: results/atm_straddle_0dte_1530/perbar_vs_pooled_qlike.csv (rv_iv notebook, section 10,
%   "Per-bar against pooled"); the FOMC split of the 14:30 bar from that cell's saved output.
\begin{table}[H]
\centering
\caption{Per-bar forecast against its pooled twin: the percent difference in
mean QLIKE (negative: the per-bar forecast is better), with a 95\% circular
block-bootstrap interval (blocks of 21 days), over the session and on three
bars labelled by bar end. DM $t$ is the Diebold--Mariano statistic of the
session mean. Bars better (worse) count the bars whose interval lies wholly
below (above) zero.}
\label{tab:app_run_perbar}
\runappfit{\small% AUTO-GENERATED by writeup/make_appendix_running_tex.py -- do not edit.
% Source: results/atm_straddle_0dte_1530/perbar_vs_pooled_qlike.csv
\begin{tabular}{lrrrrrrr}
\toprule
per-bar forecast & \shortstack{session mean\\(\%) [95\%]} & DM $t$ & \shortstack{14:30 bar\\(\%) [95\%]} & \shortstack{15:00 bar\\(\%) [95\%]} & \shortstack{16:00 bar\\(\%) [95\%]} & \shortstack{bars\\better} & \shortstack{bars\\worse} \\
\midrule
ridge (HAR + calendar) & \ensuremath{-}8.5 [\ensuremath{-}11.0, \ensuremath{-}6.0] & \ensuremath{-}5.91 & \ensuremath{-}30.5 [\ensuremath{-}34.4, \ensuremath{-}26.1] & \ensuremath{-}10.4 [\ensuremath{-}14.6, \ensuremath{-}5.8] & \ensuremath{-}6.3 [\ensuremath{-}10.4, \ensuremath{-}1.9] & 3 & 2 \\
ridge (all features) & \ensuremath{-}4.1 [\ensuremath{-}9.2, 0.8] & \ensuremath{-}1.59 & \ensuremath{-}35.7 [\ensuremath{-}47.1, \ensuremath{-}23.5] & \ensuremath{-}7.9 [\ensuremath{-}15.3, 0.7] & \ensuremath{-}2.6 [\ensuremath{-}9.4, 4.4] & 1 & 5 \\
lasso (all features) & \ensuremath{-}8.8 [\ensuremath{-}14.1, \ensuremath{-}4.1] & \ensuremath{-}3.30 & \ensuremath{-}36.1 [\ensuremath{-}49.7, \ensuremath{-}21.7] & \ensuremath{-}12.7 [\ensuremath{-}21.4, \ensuremath{-}2.7] & \ensuremath{-}2.5 [\ensuremath{-}6.7, 1.7] & 2 & 1 \\
elastic net (all features) & \ensuremath{-}9.0 [\ensuremath{-}14.2, \ensuremath{-}4.3] & \ensuremath{-}3.28 & \ensuremath{-}38.0 [\ensuremath{-}51.2, \ensuremath{-}24.3] & \ensuremath{-}14.8 [\ensuremath{-}22.5, \ensuremath{-}6.6] & \ensuremath{-}0.7 [\ensuremath{-}6.4, 5.9] & 3 & 1 \\
ridge (live-feasible set) & \ensuremath{-}3.8 [\ensuremath{-}9.6, 2.2] & \ensuremath{-}1.28 & \ensuremath{-}37.9 [\ensuremath{-}49.4, \ensuremath{-}26.1] & \ensuremath{-}8.9 [\ensuremath{-}16.1, \ensuremath{-}1.0] & \ensuremath{+}1.6 [\ensuremath{-}3.9, 7.9] & 2 & 6 \\
lasso (live-feasible set) & \ensuremath{-}3.3 [\ensuremath{-}11.0, 4.3] & \ensuremath{-}0.88 & \ensuremath{-}38.9 [\ensuremath{-}51.0, \ensuremath{-}26.4] & \ensuremath{-}9.7 [\ensuremath{-}18.4, 0.5] & \ensuremath{+}2.7 [\ensuremath{-}3.7, 10.6] & 1 & 4 \\
elastic net (live-feasible set) & \ensuremath{-}4.4 [\ensuremath{-}11.3, 2.5] & \ensuremath{-}1.29 & \ensuremath{-}39.0 [\ensuremath{-}50.9, \ensuremath{-}26.2] & \ensuremath{-}12.9 [\ensuremath{-}20.9, \ensuremath{-}3.9] & \ensuremath{+}0.9 [\ensuremath{-}3.8, 6.7] & 2 & 6 \\
\bottomrule
\end{tabular}
}
\end{table}

\paragraph{Finding.} Over the session every per-bar forecast has the lower
loss (\runappPbSessLo{}\% to \runappPbSessHi{}\%), with the interval below zero
for \runappPbSessCi{} of the \runappPbN{}. The gain sits mostly in one bar,
14:00--14:30 (\runappPbFomcBarLo{}\% to \runappPbFomcBarHi{}\%), the bar that
holds the FOMC statement: on the \runappPbFomcDays{} statement days the
difference there is \runappPbFomcLo{}\% to \runappPbFomcHi{}\%, those days
carry \runappPbShareLo{} to \runappPbShareHi{} of the pooled forecasts' loss at
that bar, and on the other \runappPbOtherDays{} days it is \runappPbOtherLo{}\%
to \runappPbOtherHi{}\%. On the traded 16:00 bar only the
\runappPbSixteenWho{} resolves a gain (\runappPbSixteen{}\%,
\runappPbSixteenCi{}).

\paragraph{Why it is not in the main text.} The per-bar gain is mostly a
calendar effect at a bar the trade does not use.

% ----------------------------------------------------------------------------
\subsection{Remaining-Variance Forecasts Across the Session}
\label{sec:app_run_multihorizon}

\paragraph{Question.} How well does each one-bar model extend to the variance
remaining until the close, at every entry clock, using only what is known at
the clock? Three regressions for the log of the remaining variance, M1 to M3
of Table~\ref{tab:app_run_multihorizon}, are fitted per entry clock on
strictly earlier days (at least 63) and rescaled by the running mean of
realized over forecast; M0 uses no model. Each is built for the HAR OLS
baseline and for the block-diagonal ridge, on the full panel: \runappMhDays{}
days from \runappMhFirst{} to \runappMhLast{}.

% Source: results/multihorizon/ttc_close_fomc1_pooled.csv and ttc_close_fomc1_byhour.csv, the
%   FIXED variant (experiments/multihorizon_ttc.py --through-close --blk2-file yhat_blk2_fomc1.parquet
%   --tag _close_fomc1), day ranges from ttc_close_fomc1_forecasts.parquet.
% The plain variant (results/multihorizon/ttc_{pooled,byhour}.csv, and the tables that
%   experiments/multihorizon_tex.py builds from it) has two defects and is not used: (1) its
%   rest-of-day target stops at the bar ending 15:30 and so omits the 15:30-16:00 bar (a median
%   9.4% of the rest of the day at 10:00 and 54.6% at 15:00, saved output of the holdclose
%   notebook, section 8d; also the macros runappMhMissTen / runappMhMissFifteen); (2) its "blk2"
%   is the ridge fitted without the FOMC calendar columns, not the panel of record.
\begin{table}[H]
\centering
\caption{QLIKE of the rest-of-day forecasts, pooled over the entry clocks
10:00--15:30 (the target runs to the 16:00 close), on all
\runappMhDays{} days and on the \runappMhPostDays{} days from 2016. DM $t$
compares the block-diagonal ridge with the HAR OLS baseline in the same
construction (negative: the ridge is better).}
\label{tab:app_run_multihorizon}
\runappfit{\small% AUTO-GENERATED by writeup/make_appendix_running_tex.py -- do not edit.
% Source: results/multihorizon/ttc_close_fomc1_pooled.csv (the fixed variant: --through-close, FOMC panel of record)
% Source: results/multihorizon/ttc_close_fomc1_forecasts.parquet (day ranges)
\begin{tabular}{lrrrrrr}
\toprule
construction & \shortstack{all days\\HAR OLS} & \shortstack{all days\\block ridge} & \shortstack{all days\\DM $t$} & \shortstack{from 2016\\HAR OLS} & \shortstack{from 2016\\block ridge} & \shortstack{from 2016\\DM $t$} \\
\midrule
M0: per-clock mean of past $\log RV^{\mathrm{rem}}$ (no model) & 0.8626 & --- & --- & 0.8993 & --- & --- \\
M1: $\log$ of the next-bar forecast & 0.1600 & 0.1435 & \ensuremath{-}5.23 & 0.2137 & 0.1877 & \ensuremath{-}3.24 \\
M2: M1 + realized variance so far & 0.1563 & 0.1407 & \ensuremath{-}6.68 & 0.2106 & 0.1858 & \ensuremath{-}4.17 \\
M2, one fit with clock dummies & 0.1565 & 0.1410 & \ensuremath{-}6.71 & 0.2099 & 0.1856 & \ensuremath{-}4.15 \\
M3: M2 + yesterday's $RV^{\mathrm{rem}}$ & 0.1507 & 0.1367 & \ensuremath{-}7.19 & 0.1997 & 0.1784 & \ensuremath{-}4.34 \\
\bottomrule
\end{tabular}
}
\end{table}

\paragraph{Finding.} The block-diagonal ridge beats the baseline in every
construction and era ($|t|$ from \runappMhDmLo{} to \runappMhDmHi{}); M1
already has \runappMhBOne{} against \runappMhAOne{}, and M3 is the best,
\runappMhBThree{}, against \runappMhNaive{} for the model-free clock mean. The
loss drops between the 14:00 and 14:30 entries (M1, block ridge:
\runappMhFourteen{} to \runappMhFourteenThirty{}) because the 14:00--14:30
bar, the FOMC statement bar, leaves the window. On the option days the same
maps do not trade better (Table~\ref{tab:app_run_maps}): M3, the best here,
trades worst there.

\paragraph{Why it is not in the main text.} The paper's rest-of-session
sections are parked.

% ----------------------------------------------------------------------------
\subsection{Sizing by the Size of the Forecast Gap}
\label{sec:app_run_sizing}

\paragraph{Question.} $\mathrm{sign}(s)$ bets the same amount every day. Does
a position that grows with the gap do better? Three sizes, each set from
earlier days only: (a) $s_d$ over the root mean square of $s$ on earlier days
(unbounded); (b) $\mathrm{sign}(s_d)$ times the share of earlier days with a
smaller $|s|$; (c) twice the share of earlier days with a smaller $s$, minus
one, which can take the side opposite to $\mathrm{sign}(s)$. The first block
scores the days after a year of history, the second every day from the
second.

% Source: results/atm_straddle_0dte_1530/vrp_sized_rules.csv (rv_iv notebook, section 10,
%   "Sizing by the size of the forecast gap").
\begin{table}[H]
\centering
\caption{Sized rules against $\mathrm{sign}(s)$ on the same days. The ridge
columns give the block-diagonal ridge's sized Sharpe ratio and its difference
from $\mathrm{sign}(s)$ with a 95\% circular block-bootstrap interval (blocks
of 21 days). The count columns report, over the \runappSzForecasts{}
forecasts, how many differences are positive, and how many intervals lie
wholly above and wholly below zero.}
\label{tab:app_run_sizing}
\runappfit{\small% AUTO-GENERATED by writeup/make_appendix_running_tex.py -- do not edit.
% Source: results/atm_straddle_0dte_1530/vrp_sized_rules.csv
\begin{tabular}{llrrrrrr}
\toprule
size & sample & days & \shortstack{ridge\\Sharpe (mid)} & \shortstack{ridge: minus $\mathrm{sign}(s)$\\(mid) [95\%]} & \shortstack{15 forecasts (mid):\\$>0$ / CI$>0$ / CI$<0$} & \shortstack{median\\(mid)} & \shortstack{15 forecasts (crossed):\\$>0$ / CI$>0$ / CI$<0$} \\
\midrule
(a) proportional, $s_d/\mathrm{rms}_d$ & after 252 days & 614 & 0.76 & \ensuremath{-}0.78 [\ensuremath{-}2.09, 0.62] & 0 / 0 / 0 & \ensuremath{-}0.73 & 0 / 0 / 0 \\
(b) rank of $|s_d|$, sign kept & after 252 days & 614 & 1.58 & \ensuremath{+}0.04 [\ensuremath{-}0.73, 0.86] & 8 / 0 / 0 & \ensuremath{+}0.03 & 8 / 0 / 0 \\
(c) centred rank, $2\hat F_d(s_d)-1$ & after 252 days & 614 & 1.14 & \ensuremath{-}0.39 [\ensuremath{-}1.36, 0.51] & 0 / 0 / 0 & \ensuremath{-}0.38 & 0 / 0 / 0 \\
\midrule
(a) proportional, $s_d/\mathrm{rms}_d$ & from day 2 & 865 & \ensuremath{-}0.70 & \ensuremath{-}2.05 [\ensuremath{-}3.05, \ensuremath{-}0.26] & 0 / 0 / 11 & \ensuremath{-}1.84 & 0 / 0 / 10 \\
(b) rank of $|s_d|$, sign kept & from day 2 & 865 & 1.13 & \ensuremath{-}0.23 [\ensuremath{-}0.90, 0.49] & 2 / 0 / 0 & \ensuremath{-}0.13 & 2 / 0 / 0 \\
(c) centred rank, $2\hat F_d(s_d)-1$ & from day 2 & 865 & 0.62 & \ensuremath{-}0.74 [\ensuremath{-}1.52, 0.14] & 0 / 0 / 5 & \ensuremath{-}0.54 & 0 / 0 / 3 \\
\bottomrule
\end{tabular}
}
\end{table}

\paragraph{Finding.} After a year of history, proportional sizing loses to
$\mathrm{sign}(s)$ for all \runappSzForecasts{} forecasts (block-diagonal
ridge \runappSzSignMid{} for $\mathrm{sign}(s)$ on those \runappSzDays{}
days); the rank with the sign kept ties it (\runappSzRankAbove{} of
\runappSzForecasts{} above, median \runappSzRankMed{}); the centred rank is
below it for all \runappSzForecasts{}. No sized rule has an interval above
zero for any forecast, block or fill.

\paragraph{Why it is not in the main text.} It is a negative result; the main
text states only that the position uses the sign.

% ----------------------------------------------------------------------------
\subsection{The Trade Beside the S\&P~500 and Other Markets}
\label{sec:app_run_portfolio}

\paragraph{Question.} An investor already holds the S\&P~500. How much of the
$\mathrm{sign}(s)$ trade belongs beside it, and is the trade a disguised bet
on another market? The weight is the maximum-Sharpe mix of the two
risk-scaled daily returns, re-estimated every day from earlier days only
(252-day warm-up, then expanding; no short positions), at the midpoint.

% Source: results/atm_straddle_0dte_1530/portfolio_weights_causal.csv (rv_iv notebook,
%   section 18).
\begin{table}[H]
\centering
\caption{The trade beside the S\&P~500, \runappPwDays{} days from
\runappPwFirst{}. Correlations are those of the trade's daily return with the
S\&P~500 over the trade's own window (15:30 to the close), with the size of
that move, and close to close. The risk share is the trade's share of the risk
budget; the dollar column is the median premium held per \$100 of index. The
last column is the combination's Sharpe ratio minus that of the S\&P~500 alone
(\runappPwSpx{}), with a 95\% paired block-bootstrap interval.}
\label{tab:app_run_portfolio}
\runappfit{\small% AUTO-GENERATED by writeup/make_appendix_running_tex.py -- do not edit.
% Source: results/atm_straddle_0dte_1530/portfolio_weights_causal.csv
\begin{tabular}{lrrrrrrr}
\toprule
forecast & \shortstack{corr.\ S\&P\\15:30$\to$close} & \shortstack{corr.\ $|$S\&P$|$\\15:30$\to$close} & \shortstack{corr.\ S\&P\\close-to-close} & \shortstack{mean\\risk share} & \shortstack{premium \$\\per \$100 S\&P} & \shortstack{Sharpe\\combined} & \shortstack{minus S\&P alone\\{}[95\%]} \\
\midrule
baseline (HAR + calendar OLS) & 0.04 & \ensuremath{-}0.23 & 0.00 & 0.63 & 3.42 & 1.00 & \ensuremath{+}0.60 [\ensuremath{-}0.67, 1.77] \\
block-diagonal ridge & \ensuremath{-}0.04 & \ensuremath{-}0.14 & \ensuremath{-}0.01 & 0.73 & 4.32 & 1.44 & \ensuremath{+}1.04 [\ensuremath{-}0.46, 2.55] \\
block-diagonal ridge, without the FOMC columns & \ensuremath{-}0.06 & \ensuremath{-}0.14 & \ensuremath{-}0.01 & 0.68 & 3.82 & 1.40 & \ensuremath{+}1.00 [\ensuremath{-}0.32, 2.30] \\
LightGBM & \ensuremath{-}0.07 & \ensuremath{-}0.08 & \ensuremath{-}0.01 & 0.81 & 6.82 & 1.37 & \ensuremath{+}0.96 [\ensuremath{-}0.70, 2.52] \\
XGBoost & \ensuremath{-}0.04 & \ensuremath{-}0.11 & \ensuremath{-}0.01 & 0.81 & 7.11 & 1.48 & \ensuremath{+}1.08 [\ensuremath{-}0.51, 2.57] \\
lasso (causally tuned) & \ensuremath{-}0.05 & \ensuremath{-}0.09 & \ensuremath{-}0.04 & 0.69 & 4.28 & 1.68 & \ensuremath{+}1.28 [\ensuremath{-}0.14, 2.66] \\
lasso ($\alpha=10^{-4}$) & \ensuremath{-}0.05 & \ensuremath{-}0.08 & \ensuremath{-}0.04 & 0.76 & 5.02 & 1.58 & \ensuremath{+}1.18 [\ensuremath{-}0.33, 2.66] \\
elastic net (causally tuned) & 0.00 & \ensuremath{-}0.14 & \ensuremath{-}0.03 & 0.66 & 3.22 & 1.08 & \ensuremath{+}0.67 [\ensuremath{-}0.74, 2.02] \\
per-bar ridge (HAR + calendar) & 0.08 & \ensuremath{-}0.12 & 0.02 & 0.73 & 4.52 & 1.36 & \ensuremath{+}0.95 [\ensuremath{-}0.29, 2.06] \\
per-bar ridge (all features) & \ensuremath{-}0.05 & \ensuremath{-}0.16 & \ensuremath{-}0.04 & 0.80 & 5.66 & 1.43 & \ensuremath{+}1.03 [\ensuremath{-}0.41, 2.45] \\
per-bar lasso (all features) & \ensuremath{-}0.05 & \ensuremath{-}0.21 & \ensuremath{-}0.02 & 0.79 & 5.87 & 1.34 & \ensuremath{+}0.93 [\ensuremath{-}0.57, 2.38] \\
per-bar elastic net (all features) & \ensuremath{-}0.05 & \ensuremath{-}0.19 & \ensuremath{-}0.02 & 0.81 & 5.68 & 1.28 & \ensuremath{+}0.87 [\ensuremath{-}0.66, 2.31] \\
per-bar ridge (live-feasible set) & \ensuremath{-}0.11 & \ensuremath{-}0.17 & 0.01 & 0.78 & 5.52 & 1.49 & \ensuremath{+}1.09 [\ensuremath{-}0.43, 2.54] \\
per-bar lasso (live-feasible set) & \ensuremath{-}0.06 & \ensuremath{-}0.16 & 0.03 & 0.75 & 5.12 & 1.63 & \ensuremath{+}1.23 [\ensuremath{-}0.31, 2.74] \\
per-bar elastic net (live-feasible set) & \ensuremath{-}0.08 & \ensuremath{-}0.17 & 0.01 & 0.77 & 5.38 & 1.71 & \ensuremath{+}1.31 [\ensuremath{-}0.11, 2.71] \\
\bottomrule
\end{tabular}
}
\end{table}

% Source: results/atm_straddle_0dte_1530/trade_return_comoves.csv (rv_iv notebook, section 18,
%   "What the trade's return moves with"); a change of exactly zero is a repeated print and
%   is dropped for that series.
\begin{table}[H]
\centering
\caption{Correlation of the trade's daily return (midpoint) with same-day
moves of other markets, on the \runappHcNdays{} days wherever both exist, with
a 95\% circular block-bootstrap interval (blocks of 21 days).}
\label{tab:app_run_comoves}
\runappfit{\small% AUTO-GENERATED by writeup/make_appendix_running_tex.py -- do not edit.
% Source: results/atm_straddle_0dte_1530/trade_return_comoves.csv
\begin{tabular}{lrrr}
\toprule
same-day move & days & block-diagonal ridge [95\%] & per-bar ridge, live-feasible [95\%] \\
\midrule
S\&P 500 return, 15:30$\to$close & 866 & \ensuremath{-}0.060 [\ensuremath{-}0.173, 0.019] & \ensuremath{-}0.068 [\ensuremath{-}0.181, 0.017] \\
S\&P 500 $|$return$|$, 15:30$\to$close & 866 & \ensuremath{-}0.109 [\ensuremath{-}0.206, \ensuremath{-}0.001] & \ensuremath{-}0.091 [\ensuremath{-}0.200, 0.018] \\
S\&P 500 return, close-to-close & 866 & 0.001 [\ensuremath{-}0.059, 0.050] & \ensuremath{-}0.001 [\ensuremath{-}0.058, 0.066] \\
S\&P 500 $|$return$|$, close-to-close & 866 & \ensuremath{-}0.046 [\ensuremath{-}0.088, 0.010] & 0.001 [\ensuremath{-}0.046, 0.059] \\
VIX change, close-to-close (points) & 861 & 0.012 [\ensuremath{-}0.038, 0.067] & 0.009 [\ensuremath{-}0.056, 0.055] \\
VIX change, 15:30$\to$16:00 (panel) & 795 & 0.062 [\ensuremath{-}0.027, 0.155] & 0.040 [\ensuremath{-}0.038, 0.143] \\
Nasdaq-100 return, close-to-close & 866 & \ensuremath{-}0.019 [\ensuremath{-}0.079, 0.033] & \ensuremath{-}0.010 [\ensuremath{-}0.073, 0.057] \\
Russell 2000 return, close-to-close & 866 & 0.005 [\ensuremath{-}0.059, 0.059] & 0.022 [\ensuremath{-}0.043, 0.093] \\
10-year Treasury yield change, close-to-close (points) & 856 & 0.027 [\ensuremath{-}0.035, 0.086] & \ensuremath{-}0.017 [\ensuremath{-}0.086, 0.051] \\
long Treasuries (TLT) total return, close-to-close & 864 & \ensuremath{-}0.031 [\ensuremath{-}0.093, 0.035] & 0.027 [\ensuremath{-}0.036, 0.089] \\
high-yield credit (HYG) total return, close-to-close & 858 & \ensuremath{-}0.017 [\ensuremath{-}0.059, 0.034] & \ensuremath{-}0.002 [\ensuremath{-}0.062, 0.073] \\
dollar index return, close-to-close & 856 & \ensuremath{-}0.041 [\ensuremath{-}0.100, 0.016] & \ensuremath{-}0.043 [\ensuremath{-}0.112, 0.023] \\
realized minus implied variance, 15:30$\to$16:00 & 866 & \ensuremath{-}0.013 [\ensuremath{-}0.065, 0.101] & \ensuremath{-}0.030 [\ensuremath{-}0.084, 0.090] \\
\bottomrule
\end{tabular}
}
\end{table}

\paragraph{Finding.} The trade barely moves with the index's direction
(close-to-close correlation \runappPwCorrLo{} to \runappPwCorrHi{} across the
\runappPwForecasts{} forecasts) and moves against the size of the last half
hour's move (\runappPwAbsLo{} to \runappPwAbsHi{}), since the rule sells on
most days. Earlier days put most of the risk budget in it: for the
block-diagonal ridge a mean risk share of \runappPwShare{}, a median
\$\runappPwDollar{} of premium per \$100 of index, and a combined Sharpe ratio
of \runappPwComb{} against \runappPwSpx{} --- a gain of \runappPwGain{}
\runappPwGainCi{}. Every forecast's interval includes zero (highest lower
bound \runappPwBestLo{}). Across the \runappCoSeries{} same-day series only
\runappCoExA{} interval excludes zero for the block-diagonal ridge, the size
of the S\&P~500's move from 15:30 to the close (\runappCoOne{},
\runappCoOneCi{}), and \runappCoExB{} for the per-bar ridge.

\paragraph{Why it is not in the main text.} The paper is about the forecast,
not portfolio construction, and on this sample the gain is not distinguishable
from none.

% ----------------------------------------------------------------------------
% Strategy variations (checklist item F1, another workstream): that workstream writes
% generated/strategy_variations.tex, which must carry its own \subsection heading.
\IfFileExists{generated/strategy_variations.tex}{% AUTO-GENERATED by experiments/strategy_variations_1530.py -- do not edit.
% Sources: results/strategy_variations/VARIATIONS.md (definitions), headline_per_straddle_premium.csv,
% cost_line.csv. Every number below is read from those files by the generator.
\subsection{Strategy variations of the last-30-min trade}
\label{sec:app_strategy_variations}

The straddle --- the nearest out-of-the-money call and the nearest out-of-the-money put,
same-day expiry, one position --- is one choice of instrument for the $\mathrm{sign}(s)$
rule. Table~\ref{tab:app_strategy_variations_defs} writes down the standard alternatives
and the way the rule maps onto each: a variation is written as the structure $L$ bought
when $s>0$, and it is sold (the position is $-L$) when $s<0$; always short sells it every
day. The straddle with a wing hedge is the one exception: it holds the plain straddle on
buying days and the short iron butterfly on selling days. All structures are same-day
expiry SPXW options entered at the 15:30 quotes and cash-settled at the official close,
on the same 866 days as the straddle. Strikes are the live 15:30 strikes ($K_c$ the
smallest call strike at or above the index $S$, $K_p$ the largest put strike at or below
it); a wing is the nearest live strike at least $w\in\{10,25,50\}$ index points beyond the
body strike. Fills are the midpoint of every leg, or crossed (every leg bought at the ask
and sold at the bid). The return is the day's profit and loss in index points divided by
that day's straddle premium at the same fill, so every row shares the straddle row's
denominator and a paired difference isolates the structure; the straddle row is the
straddle trade's return. Intervals are 95\% circular block bootstrap intervals (block 21
days, 2{,}000 draws), paired against the straddle on the same days, rule and fill.

\begingroup\footnotesize\setlength{\tabcolsep}{3pt}
\begin{longtable}{p{2.7cm}p{4.3cm}p{3.4cm}p{1.6cm}p{3.4cm}}
\caption{The strategy variations. $C(K)$, $P(K)$: a call, a put at strike $K$; $+$ bought,
$-$ sold in the buy-side form $L$. $K_c^{+}$, $K_p^{-}$: the next live strike beyond $K_c$,
$K_p$.}\label{tab:app_strategy_variations_defs}\\
\toprule
variation & bought when $s>0$ ($L$) & sold when $s<0$ ($-L$) & own notional & purpose \\
\midrule
\endfirsthead
\caption[]{(continued)}\\
\toprule
variation & bought when $s>0$ ($L$) & sold when $s<0$ ($-L$) & own notional & purpose \\
\midrule
\endhead
straddle (baseline) & $+C(K_c)+P(K_p)$ & short straddle & premium & convexity both ways at the money \\
strangle, one strike wide & $+C(K_c^{+})+P(K_p^{-})$ & short strangle & premium & all time value; cheaper, less gamma at the money \\
strangle $x$ OTM, $x\in\{0.5,1,2\}\%$ & $+C(\ge S(1+x))+P(\le S(1-x))$ & short strangle & premium & tail-only exposure; sold, a small premium for the tails \\
iron butterfly, $w$ & $+C(K_c)+P(K_p)-C(\ge K_c+w)-P(\le K_p-w)$ & short straddle + long wings (credit) & larger wing gap & tail cap on the selling side \\
iron condor, $w$ & body: the strangle one strike wide; wings $w$ beyond it & short strangle + long wings (credit) & larger wing gap & cheaper body, tail cap \\
put vertical, $w$ & $+P(K_p)-P(\le K_p-w)$ & credit put vertical & wing gap & premium from the put side alone, capped \\
call vertical, $w$ & $+C(K_c)-C(\ge K_c+w)$ & credit call vertical & wing gap & premium from the call side alone, capped \\
call (put) butterfly, $w$ & short butterfly $-C(K_0-w)+2C(K_0)-C(K_0+w)$; $K_0$ the straddle strike nearer $S$ & long butterfly (debit; pays on a pin at $K_0$) & $w$ & defined-risk selling with a pin payoff; one wing in the money \\
put (call) ratio $1\times2$, $w$ & backspread $-P(K_p)+2P(\le K_p-w)$ & front ratio $+P(K_p)-2P(\le K_p-w)$, one leg uncovered & none & tail convexity when bought, tail premium when sold \\
straddle + wing hedge, $w$ & the straddle & the short iron butterfly & none (two units) & cap the selling side's tail, keep the buying side's convexity \\
\bottomrule
\end{longtable}
\endgroup

\paragraph{What the variations show.} Under $\mathrm{sign}(s)$ with the per-bar ridge
(Table~\ref{tab:app_strategy_variations_res}), the highest Sharpe ratio at the midpoint is
the iron butterfly, w50 row (1.94; the straddle
1.90), and at crossed fills the straddle row
(1.44). The mid-fill lead of the iron butterfly, w50 row over the straddle is 0.04 [$-$0.04, 0.14]. Wings cost the rule at crossed fills: 6
of the 6 iron-butterfly and wing-hedge rows lie below the straddle with an interval
that excludes zero, against 2 of 6 at the midpoint, where the cost shrinks
as the wing moves out. The median entry cost (crossed minus mid) is
2.9\% of the straddle's own premium, 6.1\% (w10) falling to 4.0\% (w50) of the iron
butterfly's, 10.4\% falling to 7.0\% of the iron condor's, and the single-type
butterflies, with one wing in the money, cost 9.7\%--13.5\% of the straddle's
premium against 2.9\% for the straddle itself.

\begingroup\footnotesize\setlength{\tabcolsep}{1pt}
\begin{longtable}{lrrlrrlrr}
\caption{Strategy variations under $\mathrm{sign}(s)$ with the per-bar ridge on the live-feasible set (16:00-bar recalibration), 866 days (861 for the butterflies at $w=25,50$). Sharpe ratios annualized by $\sqrt{252}$ per trade day, return per straddle premium; $\Delta$Sharpe: the variation minus the straddle, paired. Worst day: mid fill, in straddle premiums. Cost share: the median crossed-minus-mid entry cost over the structure's own absolute midpoint premium, averaged over the buying and the selling side.}\label{tab:app_strategy_variations_res}\\
\toprule
& \multicolumn{3}{c}{mid fill} & \multicolumn{3}{c}{crossed fill} & worst & cost \\
\cmidrule(lr){2-4}\cmidrule(lr){5-7}
variation & Sharpe & $\Delta$Sharpe & [95\% interval] & Sharpe & $\Delta$Sharpe & [95\% interval] & day & share \\
\midrule
\endfirsthead
\caption[]{(continued)}\\
\toprule
& \multicolumn{3}{c}{mid fill} & \multicolumn{3}{c}{crossed fill} & worst & cost \\
\cmidrule(lr){2-4}\cmidrule(lr){5-7}
variation & Sharpe & $\Delta$Sharpe & [95\% interval] & Sharpe & $\Delta$Sharpe & [95\% interval] & day & share \\
\midrule
\endhead
straddle & 1.90 & --- &  & 1.44 & --- &  & $-$5.42 & 0.03 \\
\addlinespace[2pt]
strangle, one strike wide & 1.81 & $-$0.10 & [$-$0.54, 0.31] & 1.37 & $-$0.07 & [$-$0.53, 0.34] & $-$4.68 & 0.04 \\
strangle 0.5\% OTM & 1.19 & $-$0.71 & [$-$1.94, 0.52] & 0.33 & $-$1.11 & [$-$2.37, 0.09] & $-$2.87 & 0.25 \\
strangle 1\% OTM & 0.04 & $-$1.87 & [$-$3.15, 0.38] & $-$1.42 & $-$2.86 & [$-$5.41, $-$1.67] & $-$1.77 & 0.67 \\
strangle 2\% OTM & $-$0.57 & $-$2.47 & [$-$3.69, 0.06] & $-$3.01 & $-$4.45 & [$-$13.44, $-$3.16] & $-$1.43 & 1.00 \\
\addlinespace[2pt]
iron butterfly, w10 & 1.66 & $-$0.24 & [$-$0.73, 0.23] & 0.58 & $-$0.86 & [$-$1.39, $-$0.37] & $-$2.43 & 0.06 \\
iron butterfly, w25 & 1.84 & $-$0.06 & [$-$0.23, 0.12] & 1.13 & $-$0.31 & [$-$0.48, $-$0.13] & $-$5.44 & 0.04 \\
iron butterfly, w50 & 1.94 & 0.04 & [$-$0.04, 0.14] & 1.31 & $-$0.13 & [$-$0.21, $-$0.03] & $-$5.44 & 0.04 \\
\addlinespace[2pt]
iron condor, w10 & 1.69 & $-$0.21 & [$-$0.67, 0.21] & 0.56 & $-$0.88 & [$-$1.36, $-$0.46] & $-$2.53 & 0.10 \\
iron condor, w25 & 1.82 & $-$0.09 & [$-$0.53, 0.36] & 1.06 & $-$0.38 & [$-$0.84, 0.06] & $-$4.70 & 0.08 \\
iron condor, w50 & 1.85 & $-$0.05 & [$-$0.49, 0.34] & 1.18 & $-$0.26 & [$-$0.72, 0.14] & $-$4.70 & 0.07 \\
\addlinespace[2pt]
put vertical, w10 & 1.09 & $-$0.82 & [$-$1.71, 0.10] & 0.44 & $-$1.00 & [$-$1.93, $-$0.05] & $-$2.96 & 0.07 \\
put vertical, w25 & 1.31 & $-$0.59 & [$-$1.36, 0.19] & 0.88 & $-$0.56 & [$-$1.35, 0.24] & $-$4.03 & 0.05 \\
put vertical, w50 & 1.36 & $-$0.54 & [$-$1.29, 0.21] & 0.98 & $-$0.46 & [$-$1.23, 0.32] & $-$3.98 & 0.04 \\
\addlinespace[2pt]
call vertical, w10 & 0.83 & $-$1.08 & [$-$2.00, $-$0.13] & 0.22 & $-$1.22 & [$-$2.16, $-$0.23] & $-$2.58 & 0.06 \\
call vertical, w25 & 0.87 & $-$1.03 & [$-$1.88, $-$0.24] & 0.45 & $-$0.99 & [$-$1.86, $-$0.17] & $-$5.71 & 0.04 \\
call vertical, w50 & 0.96 & $-$0.94 & [$-$1.78, $-$0.17] & 0.57 & $-$0.87 & [$-$1.73, $-$0.08] & $-$5.71 & 0.04 \\
\addlinespace[2pt]
call butterfly, w10 & 1.07 & $-$0.83 & [$-$1.71, 0.01] & $-$1.33 & $-$2.77 & [$-$3.72, $-$1.85] & $-$3.84 & 0.13 \\
call butterfly, w25 & 1.48 & $-$0.37 & [$-$0.97, 0.18] & $-$0.68 & $-$2.07 & [$-$2.92, $-$1.37] & $-$4.92 & 0.04 \\
call butterfly, w50 & 1.68 & $-$0.27 & [$-$0.80, 0.21] & $-$0.68 & $-$2.17 & [$-$2.99, $-$1.46] & $-$4.92 & 0.02 \\
\addlinespace[2pt]
put butterfly, w10 & 1.05 & $-$0.85 & [$-$1.73, $-$0.01] & $-$1.50 & $-$2.94 & [$-$3.92, $-$2.01] & $-$3.84 & 0.14 \\
put butterfly, w25 & 1.51 & $-$0.35 & [$-$0.93, 0.19] & $-$0.81 & $-$2.20 & [$-$3.05, $-$1.45] & $-$4.91 & 0.04 \\
put butterfly, w50 & 1.74 & $-$0.22 & [$-$0.74, 0.27] & $-$1.13 & $-$2.62 & [$-$3.45, $-$1.90] & $-$4.89 & 0.02 \\
\addlinespace[2pt]
put ratio $1\times2$, w10 & $-$0.28 & $-$2.18 & [$-$3.82, $-$0.53] & $-$1.13 & $-$2.57 & [$-$4.21, $-$0.92] & $-$4.79 & 0.14 \\
put ratio $1\times2$, w25 & $-$1.22 & $-$3.12 & [$-$4.92, $-$1.16] & $-$1.77 & $-$3.21 & [$-$5.02, $-$1.27] & $-$7.21 & 0.06 \\
put ratio $1\times2$, w50 & $-$1.37 & $-$3.27 & [$-$5.01, $-$1.40] & $-$1.84 & $-$3.28 & [$-$5.03, $-$1.42] & $-$10.72 & 0.05 \\
\addlinespace[2pt]
call ratio $1\times2$, w10 & $-$0.43 & $-$2.34 & [$-$4.06, $-$0.66] & $-$1.34 & $-$2.78 & [$-$4.47, $-$1.15] & $-$4.60 & 0.11 \\
call ratio $1\times2$, w25 & $-$0.79 & $-$2.69 & [$-$4.48, $-$0.84] & $-$1.38 & $-$2.81 & [$-$4.62, $-$0.97] & $-$4.69 & 0.06 \\
call ratio $1\times2$, w50 & $-$0.99 & $-$2.89 & [$-$4.70, $-$1.06] & $-$1.51 & $-$2.95 & [$-$4.80, $-$1.10] & $-$5.14 & 0.05 \\
\addlinespace[2pt]
straddle + wing hedge, w10 & 1.44 & $-$0.46 & [$-$0.85, $-$0.10] & 0.73 & $-$0.71 & [$-$1.12, $-$0.36] & $-$2.43 & 0.04 \\
straddle + wing hedge, w25 & 1.69 & $-$0.21 & [$-$0.38, $-$0.01] & 1.10 & $-$0.34 & [$-$0.51, $-$0.12] & $-$5.44 & 0.04 \\
straddle + wing hedge, w50 & 1.85 & $-$0.05 & [$-$0.12, 0.06] & 1.30 & $-$0.14 & [$-$0.22, $-$0.03] & $-$5.44 & 0.03 \\
\bottomrule
\end{longtable}
\endgroup

\begingroup\footnotesize\setlength{\tabcolsep}{1pt}
\begin{longtable}{lrrlrrlrr}
\caption{As Table~\ref{tab:app_strategy_variations_res}, with $\mathrm{sign}(s)$ on the block-diagonal ridge (the paper's forecast, with the paper's recalibration over the session bars).}\label{tab:app_strategy_variations_res_blk}\\
\toprule
& \multicolumn{3}{c}{mid fill} & \multicolumn{3}{c}{crossed fill} & worst & cost \\
\cmidrule(lr){2-4}\cmidrule(lr){5-7}
variation & Sharpe & $\Delta$Sharpe & [95\% interval] & Sharpe & $\Delta$Sharpe & [95\% interval] & day & share \\
\midrule
\endfirsthead
\caption[]{(continued)}\\
\toprule
& \multicolumn{3}{c}{mid fill} & \multicolumn{3}{c}{crossed fill} & worst & cost \\
\cmidrule(lr){2-4}\cmidrule(lr){5-7}
variation & Sharpe & $\Delta$Sharpe & [95\% interval] & Sharpe & $\Delta$Sharpe & [95\% interval] & day & share \\
\midrule
\endhead
straddle & 1.34 & --- &  & 0.87 & --- &  & $-$5.42 & 0.03 \\
\addlinespace[2pt]
strangle, one strike wide & 1.31 & $-$0.03 & [$-$0.47, 0.40] & 0.86 & $-$0.01 & [$-$0.45, 0.42] & $-$4.68 & 0.04 \\
strangle 0.5\% OTM & 0.66 & $-$0.67 & [$-$1.84, 0.58] & $-$0.19 & $-$1.06 & [$-$2.25, 0.13] & $-$2.87 & 0.25 \\
strangle 1\% OTM & $-$0.36 & $-$1.70 & [$-$2.94, 0.63] & $-$1.78 & $-$2.64 & [$-$4.96, $-$1.45] & $-$1.77 & 0.67 \\
strangle 2\% OTM & $-$0.47 & $-$1.81 & [$-$2.96, 0.79] & $-$2.93 & $-$3.80 & [$-$12.80, $-$2.42] & $-$1.43 & 1.00 \\
\addlinespace[2pt]
iron butterfly, w10 & 1.23 & $-$0.11 & [$-$0.61, 0.41] & 0.15 & $-$0.72 & [$-$1.25, $-$0.21] & $-$3.88 & 0.06 \\
iron butterfly, w25 & 1.32 & $-$0.02 & [$-$0.21, 0.23] & 0.60 & $-$0.27 & [$-$0.46, $-$0.01] & $-$5.44 & 0.04 \\
iron butterfly, w50 & 1.37 & 0.03 & [$-$0.04, 0.12] & 0.73 & $-$0.14 & [$-$0.21, $-$0.05] & $-$5.44 & 0.04 \\
\addlinespace[2pt]
iron condor, w10 & 1.45 & 0.11 & [$-$0.43, 0.64] & 0.32 & $-$0.55 & [$-$1.10, $-$0.01] & $-$2.53 & 0.10 \\
iron condor, w25 & 1.35 & 0.01 & [$-$0.46, 0.47] & 0.59 & $-$0.28 & [$-$0.75, 0.16] & $-$4.70 & 0.08 \\
iron condor, w50 & 1.34 & 0.00 & [$-$0.43, 0.44] & 0.67 & $-$0.20 & [$-$0.65, 0.22] & $-$4.70 & 0.07 \\
\addlinespace[2pt]
put vertical, w10 & 1.47 & 0.13 & [$-$0.71, 0.98] & 0.83 & $-$0.04 & [$-$0.91, 0.85] & $-$2.96 & 0.07 \\
put vertical, w25 & 1.51 & 0.17 & [$-$0.56, 0.89] & 1.08 & 0.21 & [$-$0.53, 0.96] & $-$4.35 & 0.05 \\
put vertical, w50 & 1.51 & 0.18 & [$-$0.57, 0.90] & 1.14 & 0.27 & [$-$0.48, 1.01] & $-$4.33 & 0.04 \\
\addlinespace[2pt]
call vertical, w10 & $-$0.08 & $-$1.42 & [$-$2.37, $-$0.54] & $-$0.68 & $-$1.55 & [$-$2.51, $-$0.66] & $-$4.63 & 0.06 \\
call vertical, w25 & 0.02 & $-$1.32 & [$-$2.27, $-$0.40] & $-$0.40 & $-$1.26 & [$-$2.21, $-$0.35] & $-$5.71 & 0.04 \\
call vertical, w50 & 0.07 & $-$1.27 & [$-$2.21, $-$0.37] & $-$0.32 & $-$1.19 & [$-$2.13, $-$0.29] & $-$5.71 & 0.04 \\
\addlinespace[2pt]
call butterfly, w10 & 1.11 & $-$0.23 & [$-$1.06, 0.57] & $-$1.29 & $-$2.16 & [$-$3.04, $-$1.36] & $-$3.84 & 0.13 \\
call butterfly, w25 & 1.23 & $-$0.05 & [$-$0.55, 0.45] & $-$0.92 & $-$1.73 & [$-$2.37, $-$1.14] & $-$5.37 & 0.04 \\
call butterfly, w50 & 1.38 & $-$0.03 & [$-$0.48, 0.39] & $-$0.97 & $-$1.91 & [$-$2.51, $-$1.36] & $-$5.32 & 0.02 \\
\addlinespace[2pt]
put butterfly, w10 & 1.09 & $-$0.25 & [$-$1.07, 0.56] & $-$1.46 & $-$2.32 & [$-$3.23, $-$1.50] & $-$3.84 & 0.14 \\
put butterfly, w25 & 1.25 & $-$0.03 & [$-$0.54, 0.50] & $-$1.04 & $-$1.86 & [$-$2.51, $-$1.24] & $-$5.68 & 0.04 \\
put butterfly, w50 & 1.39 & $-$0.02 & [$-$0.48, 0.43] & $-$1.46 & $-$2.40 & [$-$3.10, $-$1.76] & $-$5.68 & 0.02 \\
\addlinespace[2pt]
put ratio $1\times2$, w10 & $-$0.86 & $-$2.20 & [$-$3.81, $-$0.62] & $-$1.73 & $-$2.60 & [$-$4.22, $-$1.02] & $-$4.79 & 0.14 \\
put ratio $1\times2$, w25 & $-$1.46 & $-$2.80 & [$-$4.70, $-$0.97] & $-$2.02 & $-$2.89 & [$-$4.81, $-$1.05] & $-$7.21 & 0.06 \\
put ratio $1\times2$, w50 & $-$1.53 & $-$2.86 & [$-$4.73, $-$1.02] & $-$2.01 & $-$2.88 & [$-$4.73, $-$1.04] & $-$10.72 & 0.05 \\
\addlinespace[2pt]
call ratio $1\times2$, w10 & 0.24 & $-$1.09 & [$-$2.81, 0.54] & $-$0.65 & $-$1.52 & [$-$3.25, 0.11] & $-$2.80 & 0.11 \\
call ratio $1\times2$, w25 & 0.01 & $-$1.33 & [$-$3.09, 0.35] & $-$0.57 & $-$1.44 & [$-$3.25, 0.26] & $-$4.16 & 0.06 \\
call ratio $1\times2$, w50 & $-$0.09 & $-$1.43 & [$-$3.19, 0.23] & $-$0.60 & $-$1.47 & [$-$3.29, 0.21] & $-$4.64 & 0.05 \\
\addlinespace[2pt]
straddle + wing hedge, w10 & 0.98 & $-$0.36 & [$-$0.79, 0.04] & 0.25 & $-$0.62 & [$-$1.08, $-$0.21] & $-$3.88 & 0.04 \\
straddle + wing hedge, w25 & 1.16 & $-$0.18 & [$-$0.36, 0.05] & 0.57 & $-$0.30 & [$-$0.49, $-$0.07] & $-$5.44 & 0.04 \\
straddle + wing hedge, w50 & 1.28 & $-$0.06 & [$-$0.13, 0.04] & 0.72 & $-$0.15 & [$-$0.23, $-$0.05] & $-$5.44 & 0.03 \\
\bottomrule
\end{longtable}
\endgroup

\begingroup\footnotesize\setlength{\tabcolsep}{1pt}
\begin{longtable}{lrrlrrlrr}
\caption{As Table~\ref{tab:app_strategy_variations_res}, always short (no forecast): every structure sold every day.}\label{tab:app_strategy_variations_res_short}\\
\toprule
& \multicolumn{3}{c}{mid fill} & \multicolumn{3}{c}{crossed fill} & worst & cost \\
\cmidrule(lr){2-4}\cmidrule(lr){5-7}
variation & Sharpe & $\Delta$Sharpe & [95\% interval] & Sharpe & $\Delta$Sharpe & [95\% interval] & day & share \\
\midrule
\endfirsthead
\caption[]{(continued)}\\
\toprule
& \multicolumn{3}{c}{mid fill} & \multicolumn{3}{c}{crossed fill} & worst & cost \\
\cmidrule(lr){2-4}\cmidrule(lr){5-7}
variation & Sharpe & $\Delta$Sharpe & [95\% interval] & Sharpe & $\Delta$Sharpe & [95\% interval] & day & share \\
\midrule
\endhead
straddle & 0.20 & --- &  & $-$0.27 & --- &  & $-$10.32 & 0.03 \\
\addlinespace[2pt]
strangle, one strike wide & 0.63 & 0.43 & [$-$0.01, 0.93] & 0.15 & 0.43 & [0.00, 0.93] & $-$9.47 & 0.04 \\
strangle 0.5\% OTM & 1.78 & 1.57 & [0.23, 3.99] & 0.83 & 1.10 & [$-$0.10, 3.20] & $-$3.65 & 0.25 \\
strangle 1\% OTM & 1.40 & 1.20 & [$-$0.57, 16.83] & $-$0.20 & 0.08 & [$-$1.19, 7.23] & $-$1.77 & 0.67 \\
strangle 2\% OTM & 2.64 & 2.44 & [0.39, 32.95] & $-$0.55 & $-$0.28 & [$-$1.18, 4.33] & $-$1.43 & 1.00 \\
\addlinespace[2pt]
iron butterfly, w10 & $-$0.40 & $-$0.61 & [$-$1.19, $-$0.06] & $-$1.47 & $-$1.19 & [$-$1.74, $-$0.69] & $-$5.66 & 0.06 \\
iron butterfly, w25 & $-$0.11 & $-$0.31 & [$-$0.56, $-$0.04] & $-$0.82 & $-$0.55 & [$-$0.78, $-$0.29] & $-$6.87 & 0.04 \\
iron butterfly, w50 & 0.06 & $-$0.15 & [$-$0.22, $-$0.05] & $-$0.59 & $-$0.31 & [$-$0.39, $-$0.22] & $-$10.33 & 0.04 \\
\addlinespace[2pt]
iron condor, w10 & 0.07 & $-$0.14 & [$-$0.59, 0.32] & $-$1.07 & $-$0.79 & [$-$1.22, $-$0.37] & $-$4.38 & 0.10 \\
iron condor, w25 & 0.32 & 0.12 & [$-$0.30, 0.60] & $-$0.45 & $-$0.18 & [$-$0.56, 0.27] & $-$7.34 & 0.08 \\
iron condor, w50 & 0.43 & 0.23 & [$-$0.19, 0.72] & $-$0.26 & 0.01 & [$-$0.38, 0.46] & $-$9.48 & 0.07 \\
\addlinespace[2pt]
put vertical, w10 & $-$0.79 & $-$1.00 & [$-$1.99, $-$0.11] & $-$1.41 & $-$1.14 & [$-$2.10, $-$0.28] & $-$5.98 & 0.07 \\
put vertical, w25 & $-$0.48 & $-$0.69 & [$-$1.58, 0.14] & $-$0.90 & $-$0.63 & [$-$1.50, 0.19] & $-$7.20 & 0.05 \\
put vertical, w50 & $-$0.29 & $-$0.50 & [$-$1.32, 0.33] & $-$0.67 & $-$0.40 & [$-$1.21, 0.42] & $-$10.71 & 0.04 \\
\addlinespace[2pt]
call vertical, w10 & 0.35 & 0.15 & [$-$0.79, 1.05] & $-$0.27 & 0.01 & [$-$0.90, 0.87] & $-$4.63 & 0.06 \\
call vertical, w25 & 0.38 & 0.18 & [$-$0.66, 1.05] & $-$0.05 & 0.22 & [$-$0.58, 1.05] & $-$5.71 & 0.04 \\
call vertical, w50 & 0.40 & 0.19 & [$-$0.63, 1.07] & $-$0.01 & 0.27 & [$-$0.52, 1.10] & $-$5.71 & 0.04 \\
\addlinespace[2pt]
call butterfly, w10 & $-$0.29 & $-$0.50 & [$-$1.28, 0.25] & $-$2.63 & $-$2.36 & [$-$3.11, $-$1.63] & $-$4.23 & 0.13 \\
call butterfly, w25 & 0.02 & $-$0.22 & [$-$0.74, 0.28] & $-$2.12 & $-$1.88 & [$-$2.55, $-$1.29] & $-$5.95 & 0.04 \\
call butterfly, w50 & 0.24 & 0.01 & [$-$0.43, 0.42] & $-$2.07 & $-$1.83 & [$-$2.53, $-$1.24] & $-$9.61 & 0.02 \\
\addlinespace[2pt]
put butterfly, w10 & $-$0.30 & $-$0.51 & [$-$1.28, 0.22] & $-$2.79 & $-$2.51 & [$-$3.28, $-$1.78] & $-$4.31 & 0.14 \\
put butterfly, w25 & 0.02 & $-$0.22 & [$-$0.73, 0.27] & $-$2.25 & $-$2.01 & [$-$2.70, $-$1.41] & $-$5.95 & 0.04 \\
put butterfly, w50 & 0.28 & 0.04 & [$-$0.39, 0.46] & $-$2.53 & $-$2.28 & [$-$2.99, $-$1.66] & $-$9.69 & 0.02 \\
\addlinespace[2pt]
put ratio $1\times2$, w10 & 1.42 & 1.21 & [$-$0.13, 2.49] & 0.51 & 0.79 & [$-$0.59, 2.05] & $-$5.64 & 0.14 \\
put ratio $1\times2$, w25 & 0.81 & 0.61 & [$-$1.03, 2.00] & 0.24 & 0.52 & [$-$1.12, 1.89] & $-$1.22 & 0.06 \\
put ratio $1\times2$, w50 & 0.42 & 0.21 & [$-$1.46, 1.63] & $-$0.05 & 0.23 & [$-$1.47, 1.66] & $-$0.85 & 0.05 \\
\addlinespace[2pt]
call ratio $1\times2$, w10 & $-$0.10 & $-$0.31 & [$-$1.81, 1.09] & $-$1.00 & $-$0.72 & [$-$2.26, 0.69] & $-$2.49 & 0.11 \\
call ratio $1\times2$, w25 & $-$0.30 & $-$0.50 & [$-$2.22, 0.94] & $-$0.88 & $-$0.60 & [$-$2.31, 0.84] & $-$1.74 & 0.06 \\
call ratio $1\times2$, w50 & $-$0.34 & $-$0.55 & [$-$2.32, 0.96] & $-$0.85 & $-$0.58 & [$-$2.35, 0.96] & $-$0.83 & 0.05 \\
\addlinespace[2pt]
straddle + wing hedge, w10 & $-$0.40 & $-$0.61 & [$-$1.19, $-$0.06] & $-$1.47 & $-$1.19 & [$-$1.74, $-$0.69] & $-$5.66 & 0.04 \\
straddle + wing hedge, w25 & $-$0.11 & $-$0.31 & [$-$0.56, $-$0.04] & $-$0.82 & $-$0.55 & [$-$0.78, $-$0.29] & $-$6.87 & 0.04 \\
straddle + wing hedge, w50 & 0.06 & $-$0.15 & [$-$0.22, $-$0.05] & $-$0.59 & $-$0.31 & [$-$0.39, $-$0.22] & $-$10.33 & 0.03 \\
\bottomrule
\end{longtable}
\endgroup
}{}
